\chapter{Host fields, auxiliary primes and windows of scales}\label{chapter-host-fields}

Let $\psi$ be a primitive Dirichlet character. At each scale $x$ the window bound for $\psi$ is
proved inside an abelian CM field $E$, the host field, which contains the field cut out by $\psi$,
the layer of degree about $x^{2+8\varepsilon}$ of the $2$-power cyclotomic tower, and
$\mathbf Q(\sqrt{-q},\sqrt{-d})$ for two auxiliary primes $q,d$ whose logarithms are small powers of
$x$.
This chapter chooses $q$ and $d$ so that the characters of $E$ not attached to powers of $\psi$ have
square-root cancellation at the scales that matter, chooses sets of elements of the Galois group
with prescribed distributions, one of them by Behrend's construction, and shows that hosts with all
of these properties exist at every scale of a set that contains windows $[Y,Y^{3/2}]$ for
arbitrarily large $Y$.

\section{Galois data and host fields}\label{section-hf-host-fields}

\begin{definition}\label{hf-frame}
Let $e\geq0$ be an integer and put $H_e=(\mathbf Z/2^{e+2}\mathbf Z)^\times/\{\pm1\}$, a cyclic
group of order $2^e$. A \emph{Galois datum} $\Phi$ with $2$-power exponent $e$ consists of a finite abelian
group $A$ of order $n\geq1$, written additively, whose elements are numbered $0,1,\ldots,n-1$, a
surjective homomorphism $h\colon A\to H_e$, and an element $a_0\in A$ such that $h$ is injective on
the cyclic subgroup $A_0$ generated by $a_0$. In this chapter the numbering is used only to choose
base points of orbits, in Definition~\ref{hf-good-orbit}.
We write $n_0$ for the order of $a_0$ and $o=n/n_0$ for the index of $A_0$ in $A$.

The \emph{Galois group} of $\Phi$ is $G_\Phi=\{\pm1\}\times\{\pm1\}\times A$. For
$g=(\epsilon_1,\epsilon_2,a)\in G_\Phi$ we call $\epsilon_1(g)=\epsilon_1$ and
$\epsilon_2(g)=\epsilon_2$ the first and second signs of $g$. We write $\iota=(-1,-1,0)$,
$\iota_d=(1,-1,0)$, $\Psi=\{g\in G_\Phi:\epsilon_2(g)=1\}$, and $\tau_1(t)=(1,1,t)$,
$\tau_2(t)=(-1,1,t)$ for $t\in A$. Let $\theta_q\colon G_\Phi\to\{\pm1\}$ be the first sign. Let
$X^+_\Phi$ and $X^-_\Phi$ be the sets of characters $\theta\colon G_\Phi\to\mathbf C^\times$ with
$\theta(\tau_1(t))=1$ for every $t\in A_0$ and with $\theta(\iota_d)=1$ and $\theta(\iota_d)=-1$
respectively.
\end{definition}
\leanref{Formalism/SharpRestriction/RamificationFrobenius/TotientGaloisClasses/ZsmulGeneratorEqZeroOfProjection.lean}{Formalism.HC.Frame}
\leanref{Formalism/SharpRestriction/RamificationFrobenius/TotientGaloisClasses/Analytic.lean}{Formalism.HC.Frame.Xp}
\leanref{Formalism/SharpRestriction/RamificationFrobenius/TotientGaloisClasses/Analytic.lean}{Formalism.HC.Frame.Xm}

\begin{definition}\label{hf-host-field-of-character}
Let $\psi$ be a primitive Dirichlet character of conductor $N_\psi$, let $e\geq0$ be an integer, and
let $q$ and $d$ be primes. A number field $E$ is a \emph{host field of} $(\psi,e,q,d)$ if there are
an integer $N\geq1$ divisible by $2^{e+2}$, $N_\psi$, $q$ and $d$, a primitive $N$-th root of unity
$\zeta$ in $\mathbf Q(\zeta_N)$ and a field embedding $j\colon E\to\mathbf Q(\zeta_N)$ with the
following property: for every $\tau\in\operatorname{Gal}(\mathbf Q(\zeta_N)/\mathbf Q)$ and every
integer $a\geq0$ with $\tau(\zeta)=\zeta^a$, the automorphism $\tau$ fixes $j(E)$ if and only if
$a\equiv\pm1\pmod{2^{e+2}}$, $\psi(a)=1$, $\left(\frac aq\right)=1$ and
$\left(\frac ad\right)=1$.

Thus $j(E)$ is the fixed field of the subgroup of these $\tau$. When $q\equiv d\equiv3\pmod4$ it is
the compositum of $\mathbf Q(\zeta_{2^{e+2}})^+$, the field cut out by $\psi$, $\mathbf Q(\sqrt{-q})$
and $\mathbf Q(\sqrt{-d})$.
\end{definition}
\leanref{Development/AutomorphicArithmetic/CharacterBoxes/ExistsPageN.lean}{UpperTheory.Top.IsChiField}

\begin{definition}\label{hf-framed-host-field}
Let $\Phi$ be a Galois datum with $2$-power exponent $e\geq2$, group $A$ of order $n$ and Galois group
$G_\Phi$ (Definition~\ref{hf-frame}). A \emph{host field on} $\Phi$ consists of:
\begin{enumerate}
\item primes $q,d$ with $q\equiv3\pmod8$, $d\equiv7\pmod{16}$ and
  $\left(\frac{-d}{q}\right)=1$, and an integer $Q_S\geq1$ with $2^{e+2}\mid Q_S$,
  $\gcd(Q_S,qd)=1$ and $\left(\frac{-d}{\ell}\right)=1$ for every odd prime $\ell\mid Q_S$; we
  put $N=Q_Sqd$;
\item a Galois number field $E$, an isomorphism $\gamma\colon\operatorname{Gal}(E/\mathbf Q)\to
  G_\Phi$, and an element $\sqrt{-d}\in E$ with square $-d$ such that
  $\sigma(\sqrt{-d})=\epsilon_2(\gamma(\sigma))\sqrt{-d}$ for every $\sigma$;
\item a field embedding $j\colon E\to\mathbf Q(\zeta_N)$ and a homomorphism
  $\operatorname{art}\colon(\mathbf Z/N\mathbf Z)^\times\to G_\Phi$ such that
  $j\bigl(\gamma^{-1}(\operatorname{art}(u))(y)\bigr)=\sigma_u(j(y))$ for all $u$ and $y\in E$,
  where $\sigma_u$ is the automorphism $\zeta_N\mapsto\zeta_N^u$;
\item $\operatorname{art}(-1)=\iota$; $h$ maps the $A$-component of $\operatorname{art}(u)$ to the
  class of $u$ in $H_e$; $\operatorname{art}(u)\in\{1,\iota_d\}$ whenever $u\equiv1\pmod{Q_Sq}$;
  and $\operatorname{art}(u)^2=1$ whenever $u\equiv1\pmod{Q_Sd}$.
\end{enumerate}
For a character $\theta$ of $G_\Phi$ we write $\chi_\theta$ for the primitive Dirichlet character
that induces $a\mapsto\theta(\operatorname{art}(a))$ on the integers prime to $N$; these are the
\emph{characters of} $E$. We write $K=\mathbf Q(\sqrt{-d})\subset E$, $F\subset E$ for the fixed field
of $\gamma^{-1}(\iota)$, and $\bar y$ for $\gamma^{-1}(\iota)(y)$.
\end{definition}
\leanref{Formalism/SharpRestriction/RamificationFrobenius/TotientGaloisClasses/ConjEMemHostFIff.lean}{Formalism.HC.HostRaw}
\leanref{Formalism/SharpRestriction/RamificationFrobenius/TotientGaloisClasses/ConjEMemHostFIff.lean}{Formalism.HC.HostAx}
\leanref{Formalism/SharpRestriction/RamificationFrobenius/TotientGaloisClasses/Analytic.lean}{Formalism.HC.HostRaw.IsPrimOf}

\begin{proposition}\label{hf-host-field-properties}
Let $\Phi$ be a Galois datum with $2$-power exponent $e\geq2$ and group $A$ of order $n$, and let
$(q,d,Q_S,E,\gamma,j,\operatorname{art})$ be a host field on $\Phi$, with $N=Q_Sqd$. Then:
\begin{enumerate}
\item $\operatorname{art}$ is surjective, and $E$ is a CM field of degree $4n$ whose complex
  conjugation is $\gamma^{-1}(\iota)$; its maximal totally real subfield $F$ has degree $2n$;
\item $\epsilon_2(\operatorname{art}(u))=\left(\frac ud\right)$ for every $u$;
\item for every prime $p\nmid N$ and every prime $P$ of $E$ over $p$, the Frobenius of $P$ is
  $\gamma^{-1}(\operatorname{art}(p))$;
\item the inertia group of every prime of $E$ over $d$ is $\gamma^{-1}(\{1,\iota_d\})$;
\item no nonzero prime of $E$ is ramified in $E/F$;
\item there is an ideal $\mathfrak b$ of $\mathcal O_E$ with
  $\mathfrak D_{E/K}=\mathfrak b\bar{\mathfrak b}$ and $\mathfrak b+\bar{\mathfrak b}=\mathcal O_E$,
  such that $\sigma(\mathfrak b)=\mathfrak b$ whenever $\gamma(\sigma)\in\Psi$.
\end{enumerate}
\end{proposition}
\leanref{Formalism/SharpRestriction/RamificationFrobenius/TotientGaloisClasses/InertiaGaloisSymmetry/OfRawToHostRaw.lean}{Formalism.HC.HostRaw.der}
\leanref{Formalism/SharpRestriction/RamificationFrobenius/TotientGaloisClasses/ConjEMemHostFIff.lean}{Formalism.HC.HostDer}

\begin{proof}
By (iii) of Definition~\ref{hf-framed-host-field}, $\gamma^{-1}(\operatorname{art}(u))$ corresponds
under $j$ to the restriction of $\sigma_u$ to $j(E)$. Restriction from
$\operatorname{Gal}(\mathbf Q(\zeta_N)/\mathbf Q)$ to $\operatorname{Gal}(j(E)/\mathbf Q)$ is
surjective, so $\operatorname{art}$ is surjective.

For (a), $[E:\mathbf Q]=|G_\Phi|=4n$. Complex conjugation on $\mathbf Q(\zeta_N)$ is $\sigma_{-1}$, so
for every embedding of $E$ into $\mathbf C$ complex conjugation restricts to
$\gamma^{-1}(\operatorname{art}(-1))=\gamma^{-1}(\iota)$, which is not the identity. An abelian
number field on which complex conjugation acts nontrivially is a CM field, with $F$ the fixed field
of complex conjugation, so $[F:\mathbf Q]=2n$.

For (b), the square roots of $-d$ in $\mathbf Q(\zeta_N)$ are $\pm g_d$, where $g_d$ is the
quadratic Gauss sum of modulus $d$, and $\sigma_u(g_d)=\left(\frac ud\right)g_d$. Comparing with
(ii) and (iii) gives the claim.

For (c), the Frobenius of a prime of $\mathbf Q(\zeta_N)$ over $p\nmid N$ is $\sigma_p$, and its
restriction is $\gamma^{-1}(\operatorname{art}(p))$.

For (d), the inertia group at $d$ in $\mathbf Q(\zeta_N)$ consists of the $\sigma_u$ with
$u\equiv1\pmod{Q_Sq}$. Its image lies in $\gamma^{-1}(\{1,\iota_d\})$ by (iv), and it is not
trivial, since by (b) some $u\equiv1\pmod{Q_Sq}$ moves $\sqrt{-d}$.

For (e), the complex conjugation $\gamma^{-1}(\iota)$ generates $\operatorname{Gal}(E/F)$, so it
suffices that $\iota$ lies in no inertia group. At $d$ this follows from (d). At $q$ the inertia
group is the image of the $u\equiv1\pmod{Q_Sd}$, whose second sign is $\left(\frac ud\right)=1$ by
(b), while $\epsilon_2(\iota)=-1$. At a prime $\ell\mid Q_S$ the inertia group is the image of
the $u\equiv1\pmod{qd}$, whose second sign is again $1$. Primes not dividing $N$ are unramified.

For (f), the subgroup $\gamma^{-1}(\Psi)$ is $\operatorname{Gal}(E/K)$. Its intersection with the
inertia group at $d$ is trivial by (d), so every prime dividing $\mathfrak D_{E/K}$ lies over a
prime $\ell\mid Q_Sq$. Each such $\ell$ splits in $K$: for $\ell=2$ because $-d\equiv1\pmod8$, for
odd $\ell\mid Q_S$ by (i), and for $\ell=q$ because $\left(\frac{-d}{q}\right)=1$. Let
$\mathfrak b$ be the part of $\mathfrak D_{E/K}$ lying over one of the two primes of $K$ above each
such $\ell$. Then $\bar{\mathfrak b}$ is the part over the other primes, the two are coprime, and
$\operatorname{Gal}(E/K)$ fixes each prime of $K$, hence fixes $\mathfrak b$.
\end{proof}

\begin{lemma}\label{hf-kernel-host-field}
Let $\psi$ be a primitive Dirichlet character of conductor $N_\psi$. Let $\Phi$ be a Galois datum with
$2$-power exponent $e\geq2$, and let $(q,d,Q_S,E,\gamma,j,\operatorname{art})$ be a host field on
$\Phi$ (Definition~\ref{hf-framed-host-field}), with $N=Q_Sqd$. Suppose that $N_\psi\mid Q_S$ and
that for every integer $a\geq0$ prime to $N$,
\[
 \operatorname{art}(a)=1\iff a\equiv\pm1\pmod{2^{e+2}},\ \psi(a)=1,\
 \left(\tfrac aq\right)=1\text{ and }\left(\tfrac ad\right)=1.
\]
Then $E$ is a host field of $(\psi,e,q,d)$ (Definition~\ref{hf-host-field-of-character}).
\end{lemma}
\leanref{Development/AutomorphicArithmetic/CharacterBoxes/TwistForm.lean}{UpperTheory.Top.KerChi}
\leanref{Development/AutomorphicArithmetic/CharacterBoxes/IsChiField.lean}{UpperTheory.Top.isChiField}

\begin{proof}
Take $N=Q_Sqd$, which is divisible by $2^{e+2}$, $N_\psi$, $q$ and $d$, the embedding $j$, and a
primitive $N$-th root of unity $\zeta$. Let $\tau(\zeta)=\zeta^a$. Then $a$ is prime to $N$ and
$\tau=\sigma_a$. By (iii) of Definition~\ref{hf-framed-host-field} and the injectivity of $j$,
$\tau$ fixes $j(E)$ if and only if $\gamma^{-1}(\operatorname{art}(a))$ is the identity, that is,
$\operatorname{art}(a)=1$. The hypothesis translates this into the condition of
Definition~\ref{hf-host-field-of-character}.
\end{proof}

\begin{definition}\label{hf-host-field-rows}
Let $\Phi$ be a Galois datum with $2$-power exponent $e\geq2$, group of order $n$ and index $o$, let
$(q,d,Q_S,E,\gamma,j,\operatorname{art})$ be a host field on $\Phi$ with $N=Q_Sqd$, and let $B\geq0$
be an integer. The host field has \emph{splitting bound} $B$ if:
\begin{enumerate}
\item $\epsilon_1(\operatorname{art}(u))=\left(\frac uq\right)$ for every $u$, and $E$ contains an
  element $\sqrt{-q}$ with square $-q$ and $\sigma(\sqrt{-q})=\epsilon_1(\gamma(\sigma))\sqrt{-q}$
  for every $\sigma$;
\item for every prime $\ell\mid Q_S$ there is $u\equiv1\pmod{N/\ell}$ with
  $\operatorname{art}(u)\neq1$;
\item $Q_S\leq n^2$, $o^{64}\leq2^e$ and $q\neq3$;
\item $\left(\frac{-d}{p}\right)=1$ for every odd prime $p\leq B$, and $4\ell\leq B$ for every odd
  prime $\ell\mid Q_S$.
\end{enumerate}
\end{definition}
\leanref{Development/AutomorphicArithmetic/CoprimeRanks/SecondSetNoRelation.lean}{HostParam.Seam.FieldRows}

\begin{proposition}\label{hf-host-field-exists}
Let $\psi$ be a primitive Dirichlet character of conductor $N_\psi$ and order $o_\psi$. There are
integers $c\geq1$ and $e_0\geq0$ with the following property. Let $e\geq e_0$, let $q\neq3$ and $d$
be primes with $q\equiv3\pmod8$, $d\equiv7\pmod{16}$, $\left(\frac{-d}{q}\right)=1$ and
$\gcd(N_\psi,qd)=1$, and let $B\geq4N_\psi$ be an integer with $\left(\frac{-d}{p}\right)=1$ for
every odd prime $p\leq B$. Then there are a Galois datum $\Phi$ with $2$-power exponent $e$, group of order
$n=c\,2^e$ and index $o=o_\psi$, and a host field $(q,d,Q_S,E,\gamma,j,\operatorname{art})$ on
$\Phi$ with these $q$ and $d$ and with $N=Q_Sqd$, such that:
\begin{enumerate}
\item the host field has splitting bound $B$ (Definition~\ref{hf-host-field-rows});
\item $N_\psi\mid Q_S$, $Q_S\mid\operatorname{lcm}(2^{e+2},N_\psi)$ and $Q_S\leq4N_\psi n$;
\item for every integer $a\geq0$ prime to $N$, $\operatorname{art}(a)=1$ if and only if
  $a\equiv\pm1\pmod{2^{e+2}}$, $\psi(a)=1$, $\left(\frac aq\right)=1$ and
  $\left(\frac ad\right)=1$;
\item every primitive character that induces a power of $\psi$ is $\chi_\theta$ for some
  $\theta\in X^+_\Phi$, and $\chi_{\theta_q}$ induces no power of $\psi$;
\item for every character $\theta$ of $G_\Phi$ with $\theta(\tau_1(t))=1$ for all $t\in A_0$ there
  are $\xi$ in the group generated by $\psi$ and $i,j\in\{0,1\}$ with $\theta(\iota_d)=(-1)^j$ and
  $\theta(\operatorname{art}(a))=\xi(a)\left(\frac aq\right)^i\left(\frac ad\right)^j$ for every
  $a$ prime to $N$.
\end{enumerate}
\end{proposition}
\leanref{Development/AutomorphicArithmetic/CharacterBoxes/FieldChi.lean}{UpperTheory.Top.fieldChi}
\leanref{Development/AutomorphicArithmetic/CharacterBoxes/TwistForm.lean}{UpperTheory.Top.TwistForm}

\begin{proof}
Let $B_\psi$ be the group of values of $\psi$, a cyclic group of order $o_\psi$. For $e\geq0$ put
$M_e=\operatorname{lcm}(2^{e+2},N_\psi)$ and let $V_e\subset B_\psi$ be the subgroup of the values
$\psi(u)$ at integers $u$ prime to $2N_\psi$ with $u\equiv\pm1\pmod{2^{e+2}}$.

\emph{Step 1: the constants.} Since $V_{e+1}\subset V_e$, the order $|V_e|$ is a non-increasing
function of $e$ with values between $1$ and $o_\psi$. Let $c$ be its least value and $e_1$ an
integer with $|V_e|=c$ for $e\geq e_1$. Put $e_0=\max(e_1,2,4N_\psi,64o_\psi)$. Fix $e\geq e_0$
and $q,d,B$ as in the statement.

\emph{Step 2: the modulus $Q_S$.} Let $Q_S$ be the least positive integer $m$ with
$2^{e+2}\mid m\mid M_e$ such that $\psi(u)=1$ for every $u$ prime to $M_e$ with $u\equiv1\pmod m$;
the integer $m=M_e$ has these properties. Let $u$ be prime to $N_\psi$ with
$u\equiv1\pmod{\gcd(Q_S,N_\psi)}$. By the Chinese remainder theorem there is $u'$ prime to $M_e$
with $u'\equiv u\pmod{N_\psi}$ and $u'\equiv1\pmod{Q_S}$, so $\psi(u)=\psi(u')=1$. Thus $\psi$ is
induced from a character modulo $\gcd(Q_S,N_\psi)$, and since $\psi$ is primitive,
$N_\psi\mid Q_S$. Hence $\psi$ defines a homomorphism $\bar\psi\colon(\mathbf Z/Q_S\mathbf Z)^\times\to
B_\psi$, which is surjective.

\emph{Step 3: the Galois datum.} Let $\eta\colon(\mathbf Z/Q_S\mathbf Z)^\times\to H_e\times B_\psi$
send $u$ to the class of $u$ in $H_e$ and to $\bar\psi(u)$. Let $A$ be the image of $\eta$, written
additively and numbered in any way, let $h$ and $p_2$ be the two projections, and let $A_0$ be the
kernel of $p_2$. The projection $h$ is surjective and injective on $A_0$, so $A_0$ embeds in the
cyclic group $H_e$ and is cyclic; let $a_0$ be a generator. This is a Galois datum $\Phi$ with
$2$-power exponent $e$. Since $p_2$ is surjective, $o=[A:A_0]=o_\psi$. The kernel of $h$ is
$\{0\}\times V_e$, so $n=|A|=2^e|V_e|=c\,2^e$.

\emph{Step 4: the Artin map.} Put $N=Q_Sqd$. Since $Q_S\mid M_e$, $\gcd(N_\psi,qd)=1$ and $q,d$
are odd, $\gcd(Q_S,qd)=1$. Let $w\colon\{\pm1\}\to A$ send $\pm1$ to $\eta(\pm1)$; then
$2w(-1)=0$. Define $\operatorname{art}\colon(\mathbf Z/N\mathbf Z)^\times\to G_\Phi$ by
\[
 \operatorname{art}(u)=\Bigl(\bigl(\tfrac uq\bigr),\bigl(\tfrac ud\bigr),\
 \eta(u\bmod Q_S)+w\bigl(\bigl(\tfrac uq\bigr)\bigr)\Bigr).
\]
It is a homomorphism. It is surjective: given $(\epsilon_1,\epsilon_2,a)$, the Chinese remainder
theorem gives $u$ with $\eta(u\bmod Q_S)=a-w(\epsilon_1)$, $\left(\frac uq\right)=\epsilon_1$ and
$\left(\frac ud\right)=\epsilon_2$. Since $q\equiv d\equiv3\pmod4$,
$\operatorname{art}(-1)=(-1,-1,2\eta(-1))=\iota$.

\emph{Step 5: the field.} Let $E\subset\mathbf Q(\zeta_N)$ be the fixed field of the $\sigma_u$
with $\operatorname{art}(u)=1$, and let $j$ be the inclusion. Since $\operatorname{art}$ is
surjective, $\sigma_u|_E\mapsto\operatorname{art}(u)$ is an isomorphism
$\operatorname{Gal}(E/\mathbf Q)\to G_\Phi$; let $\gamma$ be this isomorphism, so that (iii) of
Definition~\ref{hf-framed-host-field} holds. Let $g_q$ and $g_d$ be the quadratic Gauss sums of
moduli $q$ and $d$. Then $g_q^2=-q$, $g_d^2=-d$, $\sigma_u(g_q)=\left(\frac uq\right)g_q$ and
$\sigma_u(g_d)=\left(\frac ud\right)g_d$. So $g_q,g_d\in E$, and with $\sqrt{-q}=g_q$ and
$\sqrt{-d}=g_d$ we get $\sigma(\sqrt{-d})=\epsilon_2(\gamma(\sigma))\sqrt{-d}$ and
$\sigma(\sqrt{-q})=\epsilon_1(\gamma(\sigma))\sqrt{-q}$ for every $\sigma$.

\emph{Step 6: the host field.} Condition (i) of Definition~\ref{hf-framed-host-field} holds by
hypothesis and Step 4; an odd prime $\ell\mid Q_S$ divides $N_\psi$, so $\ell\leq B$ and
$\left(\frac{-d}{\ell}\right)=1$. For (iv): $h$ maps the $A$-component of $\operatorname{art}(u)$ to
the class of $u$ in $H_e$, because the class of $-1$ in $H_e$ is trivial. If $u\equiv1\pmod{Q_Sq}$,
then $\operatorname{art}(u)=(1,\left(\frac ud\right),0)\in\{1,\iota_d\}$. If
$u\equiv1\pmod{Q_Sd}$, then $\operatorname{art}(u)=(\left(\frac uq\right),1,w(\left(\frac uq\right)))$,
whose square is $1$. So $(q,d,Q_S,E,\gamma,j,\operatorname{art})$ is a host field on $\Phi$.

\emph{Step 7: the splitting bound $B$.} Condition (a) of Definition~\ref{hf-host-field-rows} is
Step 4 and Step 5. For (b), let $\ell\mid Q_S$ be prime. If $2^{e+2}\mid Q_S/\ell$, then $Q_S/\ell$
fails the condition of Step 2, so there is $v$ prime to $M_e$ with $v\equiv1\pmod{Q_S/\ell}$ and
$\psi(v)\neq1$; take $u\equiv v\pmod{Q_S}$ with $u\equiv1\pmod{qd}$, so that $u\equiv1\pmod{N/\ell}$
and $p_2$ of the $A$-component of $\operatorname{art}(u)$ is $\psi(v)\neq1$. Otherwise $\ell=2$ and
$Q_S=2^{e+2}c'$ with $c'$ odd; take $u\equiv1+2^{e+1}c'\pmod{Q_S}$ with $u\equiv1\pmod{qd}$. Then
$u\equiv1\pmod{N/2}$, and $u\equiv1+2^{e+1}\not\equiv\pm1\pmod{2^{e+2}}$ since $e\geq1$, so
$\operatorname{art}(u)\neq1$. For (c): $Q_S\leq2^{e+2}N_\psi=2^e\cdot4N_\psi\leq2^{2e}\leq n^2$
because $4N_\psi\leq e<2^e$; $o^{64}=o_\psi^{64}\leq2^{64o_\psi}\leq2^e$; and $q\neq3$. For (d),
$\left(\frac{-d}{p}\right)=1$ for odd $p\leq B$ by hypothesis, and an odd prime $\ell\mid Q_S$
divides $N_\psi$, so $4\ell\leq4N_\psi\leq B$.

\emph{Step 8: properties (b) and (c).} We have $N_\psi\mid Q_S\mid M_e$, and
$Q_S\leq2^{e+2}N_\psi=4N_\psi2^e\leq4N_\psi n$. For $a\geq0$ prime to $N$,
$\operatorname{art}(a)=1$ if and only if $\left(\frac aq\right)=\left(\frac ad\right)=1$ and
$\eta(a)=0$, that is, $a\equiv\pm1\pmod{2^{e+2}}$ and $\psi(a)=1$.

\emph{Step 9: property (d).} For a character $\chi$ of $B_\psi$ put
$\theta_\chi(\epsilon_1,\epsilon_2,a)=\chi\bigl(p_2(a+w(\epsilon_1))\bigr)$. This is a character of
$G_\Phi$; it is trivial on $\tau_1(A_0)=\{1\}\times\{1\}\times\ker p_2$ and on $\iota_d$, so
$\theta_\chi\in X^+_\Phi$. Since $2w(\epsilon)=0$,
$\theta_\chi(\operatorname{art}(u))=\chi(\bar\psi(u))=\chi(\psi(u))$. Every power of $\psi$ is
$\chi\circ\psi$ for some $\chi$, so every primitive character inducing a power of $\psi$ is
$\chi_{\theta_\chi}$ for some $\chi$. The first sign satisfies
$\theta_q(\operatorname{art}(u))=\left(\frac uq\right)$, so $\chi_{\theta_q}=\left(\frac{\cdot}{q}\right)$,
which has conductor $q$. A power of $\psi$ has conductor dividing $N_\psi$, which is prime to $q$,
so $\chi_{\theta_q}$ induces no power of $\psi$.

\emph{Step 10: property (e).} The quotient of $G_\Phi$ by $\tau_1(A_0)$ is
$\{\pm1\}\times\{\pm1\}\times A/A_0$, and $p_2$ identifies $A/A_0$ with $B_\psi$; so $G_\Phi$ has
$4o_\psi$ characters trivial on $\tau_1(A_0)$. The $4o_\psi$ characters
$\theta_\chi\epsilon_1^i\epsilon_2^j$ with $i,j\in\{0,1\}$ are trivial on $\tau_1(A_0)$ and pairwise
distinct: their values at $(1,1,a)$ determine $\chi$, and then their values at $(-1,1,0)$ and at
$\iota_d=(1,-1,0)$ determine $i$ and $j$. So every $\theta$ trivial on $\tau_1(A_0)$ is
$\theta_\chi\epsilon_1^i\epsilon_2^j$, with $\theta(\iota_d)=(-1)^j$ and
$\theta(\operatorname{art}(a))=\chi(\psi(a))\left(\frac aq\right)^i\left(\frac ad\right)^j$, and
$\xi=\chi\circ\psi$ is in the group generated by $\psi$.
\end{proof}

\section{Auxiliary primes}\label{section-hf-auxiliary-primes}

For a Dirichlet character $\chi$ and $y>0$ we use the corrected Chebyshev sums of
Definition~\ref{wl-corrected-sums}: $g_\chi(y)=\sum_{n\leq y}\Lambda(n)\chi(n)(y/n-1)+y\frac{L'}{L}(1,\chi)$
for $\chi\neq1$, and $g_1(y)=G(y)=\sum_{n\leq y}\Lambda(n)(y/n-1)-y(\log y-1-\gamma)$. A primitive
character $\chi$ \emph{induces} a function $f$ on the integers prime to $M$ if $\chi(a)=f(a)$ for
all such $a$. We write $\chi_{-q}$, $\chi_{-d}$ and $\chi_8$ for the primitive quadratic characters
$\left(\frac{-q}{\cdot}\right)$, $\left(\frac{-d}{\cdot}\right)$ and $\left(\frac{2}{\cdot}\right)$;
when $q\equiv3\pmod4$, $\chi_{-q}(a)=\left(\frac aq\right)$.

\begin{definition}\label{hf-admissible-pair}
Let $\varepsilon>0$ and $Y>1$ be real numbers, and put $L=Y^{\varepsilon/5}$. A pair $(q,d)$ of
primes is \emph{admissible at} $Y$ if $q\equiv3\pmod8$, $d\equiv7\pmod{16}$,
$L/2<\log q\leq L$, $\tfrac12L\log^2Y<\log d\leq L\log^2Y$, $q\mid d+1$, and every odd prime
$p\leq L$ divides $d+1$.

If $(q,d)$ is admissible at $Y$, then $-d\equiv1$ modulo $q$ and modulo every odd prime $p\leq L$,
so $\left(\frac{-d}{q}\right)=1$ and $\left(\frac{-d}{p}\right)=1$ for these $p$.
\end{definition}
\leanref{Development/AutomorphicArithmetic/CharacterBoxes/ExistsPageN.lean}{UpperTheory.Top.WinPair}

\begin{definition}\label{hf-good-pair}
Let $M\geq1$ be an integer and let $\varepsilon>0$ and $1<Y\leq X$ be real numbers. A pair $(q,d)$
of primes is \emph{good for} $(M,\varepsilon,Y,X)$ if:
\begin{enumerate}
\item $(q,d)$ is admissible at $Y$ (Definition~\ref{hf-admissible-pair}) and $\gcd(M,qd)=1$;
\item for every Dirichlet character $\xi$ modulo $M$ and every $(i,j)\in\{0,1\}^2$ other than
  $(0,0)$, the primitive character $\chi$ that induces
  $a\mapsto\xi(a)\left(\frac aq\right)^i\left(\frac ad\right)^j$ on the integers prime to $2Mqd$
  satisfies $|g_\chi(y)|\leq y^{1/2+\varepsilon/2}$ and
  $\left|\frac{L'}{L}(1,\chi)\right|\leq4\log y$ for every $y\in[Y,X]$;
\item for every real primitive character $\rho$ of conductor dividing $8M$ and every such $(i,j)$,
  the primitive character $\chi$ that induces
  $a\mapsto\rho(a)\left(\frac aq\right)^i\left(\frac ad\right)^j$ on the integers prime to $2Mqd$
  has $L(s,\chi)\neq0$ whenever $\operatorname{Re}s>1-1/400$ and
  $|\operatorname{Im}s|\leq8M\exp(Y^{\varepsilon/5})\exp(Y^{\varepsilon/5}\log^2Y)^{j}$.
\end{enumerate}
\end{definition}
\leanref{Development/AutomorphicArithmetic/CharacterBoxes/PairGood.lean}{UpperTheory.Top.PairGood}

\begin{theorem}\label{hf-zero-density-box}
There is a constant $C$ such that for every $Q\geq3$ the number of primitive Dirichlet characters
$\chi$ of conductor at most $Q$ for which $L(s,\chi)$ has a zero $s$ with
$\operatorname{Re}s>1-1/400$ and $|\operatorname{Im}s|\leq Q$ is at most $CQ^{1/4}$.
\end{theorem}
\leanref{Formalism/GL3MeanSquares/GenusMeanAsymptotics/JacobiCharacterSums/CountZAll.lean}{Formalism.HermitianLattices.Ext.PrimitiveCharacterCasts.ZD.zero_density_box}

The log-free zero-density estimates of Gallagher and Montgomery bound the number of zeros with
real part at least $\sigma$ and height at most $Q$ of the $L$-functions of primitive characters of
conductor at most $Q$ by $Q^{c(1-\sigma)}$, up to powers of $\log Q$ in Montgomery's form; with
$1-\sigma=1/400$ this is at most $CQ^{1/4}$ \cite{Gal70}, \cite[Chapter 12]{Mon71}.

\begin{theorem}\label{hf-zero-density-regions}
There are constants $C_1,c>0$ with the following properties.
\begin{enumerate}
\item For every $Q\geq3$, the zeros $\rho$ of the functions $L(s,\chi)$, for $\chi$ primitive of
  conductor at most $Q$, with $\operatorname{Re}\rho\geq1-1/400$ and
  $|\operatorname{Im}\rho|\leq Q$, number at most $C_1Q^{1/4}$, counted with multiplicity.
\item Let $\chi$ be a primitive Dirichlet character of conductor $N\geq3$. If $L(\rho,\chi)=0$ and
  $\operatorname{Re}\rho\geq1-c/\log\max(N,N|\operatorname{Im}\rho|)$, then $\chi$ is real, $\rho$
  is real and simple, and $\rho$ is the only zero of $L(s,\chi)$ in that region.
\item Every zero $\rho$ of $\zeta$ with $0<\operatorname{Re}\rho<1$ has
  $\operatorname{Re}\rho\leq1-c/\log(|\operatorname{Im}\rho|+3)$.
\end{enumerate}
\end{theorem}
\leanref{Development/AbelianVarietyTheory/Proofs/GlobalInterpolation/DensityAndRegions.lean}{GlobalTheory.ZerosNearOne.density_and_regions}

Part (a) follows from the zero-density estimates of \cite{Gal70} and \cite[Chapter 12]{Mon71}; parts
(b) and (c) are the classical zero-free regions of de la Vallée Poussin, Landau and Page
\cite[Chapters 13 and 14]{Dav00}. These estimates are used in Chapter~\ref{chapter-assembly}.

\begin{theorem}\label{hf-primes-progressions}
For an integer $M\geq2$, an \emph{exceptional pair} modulo $M$ is a pair $(\chi_1,\beta_1)$ of a
real nontrivial Dirichlet character $\chi_1$ modulo $M$ and a real number
$1-1/(8\log M)<\beta_1<1$ with $L(\beta_1,\chi_1)=0$. Let
$\operatorname{li}X=\int_2^X dt/\log t$, and let $\pi(X;M,b)$ be the number of primes $p\leq X$
with $p\equiv b\pmod M$. There are constants $c_0,c_1,c_3>0$ and $A_0,C_0$ with the following
properties.
\begin{enumerate}
\item Let $M$ be an integer with $8\mid M$. There are an exceptional pair $(\chi_1,\beta_1)$ modulo
  $M$, or none if no exceptional pair modulo $M$ exists, such that for every real $X\geq2$ with
  $\log X\geq C_0\log M\log(\log M+2)$ and every integer $b$ prime to $M$,
  \[
   |\pi(X;M,b)-m_M(b,X)|\leq A_0\bigl(e^{-c_3\log X/\log M}+e^{-\sqrt{c_3\log X}}\bigr)m_M(b,X),
  \]
  where $m_M(b,X)=(\operatorname{li}X-\chi_1(b)\operatorname{li}X^{\beta_1})/\varphi(M)$ if there is
  the pair and $m_M(b,X)=\operatorname{li}X/\varphi(M)$ if not.
\item Let $Z\geq3$, and let $\chi_1,\chi_2$ be real nontrivial Dirichlet characters of moduli
  $M_1,M_2\leq Z$ such that $L(\beta_i,\chi_i)=0$ with $1-c_0/\log Z<\beta_i<1$ for $i=1,2$. Then
  $\chi_1(n)=\chi_2(n)$ for every $n$ prime to $M_1M_2$.
\item If $\chi$ is a real nontrivial Dirichlet character modulo $M$ and $L(\beta,\chi)=0$ with
  $\beta$ real, then $\beta\leq1-c_1/M$.
\end{enumerate}
\end{theorem}
\leanref{Formalism/GL3MeanSquares/GenusMeanAsymptotics/JacobiCharacterSums/PNTFinal.lean}{Formalism.HermitianLattices.Ext.PrimesInProgressions.pnt_progressions_relative_weak_holds}
\leanref{Formalism/GL3MeanSquares/GenusMeanAsymptotics/JacobiCharacterSums/MSqPage.lean}{Formalism.HermitianLattices.Ext.LogarithmicIntegralBounds.page_across}
\leanref{Formalism/LandauZeroFreeBounds/SiegelWalfiszAPOfSiegelLower.lean}{Formalism.ZeroFreeDirichlet.excBound}

Part (a) is the prime number theorem for arithmetic progressions with the exceptional zero in the
form of Gallagher \cite{Gal70}, see also \cite[Chapter 20]{Dav00}. Part (b) is the theorem of
Landau and Page that at most one real character has a zero in the region $1-c_0/\log Q<\beta<1$ for
the modulus $Q=M_1M_2\leq Z^2$, and part (c) follows from the bound
$1-\beta\gg M^{-1/2}(\log M)^{-2}$ \cite[Chapter 14]{Dav00}.

\begin{lemma}\label{hf-twisted-prime-sums}
For all real $\varepsilon,A,B>0$ there is $Y_0$ with the following property. Let $Y\geq Y_0$, let
$m_0\geq1$ and $a_0$ be integers with $8\mid m_0$ and $\gcd(a_0,m_0)=1$, and let $U$ satisfy
$\frac14\log^2Y\log m_0\leq\log U$ and $Y^{\varepsilon/6}\leq\log U\leq Y^{\varepsilon/5}\log^2Y$.
Let $P$ be the set of primes $q$ with $\sqrt U<q\leq U$ and $q\equiv a_0\pmod{m_0}$. Then there is
an integer $s^*$ such that
\[
 \Bigl|\sum_{q\in P}\Bigl(\frac qs\Bigr)\Bigr|\leq Y^{-B}|P|
\]
for every odd squarefree integer $1<s\leq Y^A$ prime to $m_0$ with $s\neq s^*$.
\end{lemma}
\leanref{Development/AutomorphicArithmetic/MeanSquareBounds/twistK.lean}{QH.twistK}

\begin{proof}
Let $c_0,c_1,c_3,A_0,C_0$ be the constants of Theorem~\ref{hf-primes-progressions}, and put
$Z=m_0Y^A$. Call an odd squarefree integer $1<s\leq Y^A$ prime to $m_0$ \emph{bad} if the
exceptional pair $(\chi_1,\beta_1)$ of Theorem~\ref{hf-primes-progressions}(a) for $M=m_0s$ exists,
$1-\beta_1<c_0/\log Z$, and $\chi_1(u)=\left(\frac us\right)$ for every $u\equiv1\pmod{m_0}$ prime to
$s$.

There is at most one bad integer. Let $s_1,s_2$ be bad. Their moduli $m_0s_1,m_0s_2$ are at most
$Z$, so Theorem~\ref{hf-primes-progressions}(b) shows that the two exceptional characters agree on
the integers prime to $m_0^2s_1s_2$. Hence $\left(\frac u{s_1}\right)=\left(\frac u{s_2}\right)$ for
every $u\equiv1\pmod{m_0}$ prime to $s_1s_2$. For distinct odd squarefree $s_1,s_2$ prime to $m_0$
the Chinese remainder theorem gives such a $u$ at which the two symbols differ, so $s_1=s_2$. Let
$s^*$ be the bad integer if there is one, and $s^*=0$ otherwise.

Let $s\neq s^*$ be as in the statement, so $s$ is not bad, and put $M=m_0s$. For $Y$ large,
$\log U\geq C_0\log M\log(\log M+2)$, because $\log M\leq\log m_0+A\log Y\leq4\log U/\log^2Y+A\log Y$
and $\log U\geq Y^{\varepsilon/6}$. The primes $q\in P$ lie in the classes $b$ modulo $M$ with
$b\equiv a_0\pmod{m_0}$ and $b$ prime to $s$, and
$\sum_{q\in P}\left(\frac qs\right)=\sum_b\left(\frac bs\right)\bigl(\pi(U;M,b)-\pi(\sqrt U;M,b)\bigr)$.
Replace each $\pi(U;M,b)$ by $m_M(b,U)$, with relative error at most
$k=A_0(e^{-c_3\log U/\log M}+e^{-\sqrt{c_3\log U}})$ by Theorem~\ref{hf-primes-progressions}(a).
The terms $\operatorname{li}U/\varphi(M)$ cancel: as $b$ runs over these classes, $b$ modulo $s$ runs
over all units, and $\left(\frac{\cdot}{s}\right)$ is a nontrivial character modulo $s$. The
exceptional terms are $\left(\frac bs\right)\chi_1(b)\operatorname{li}U^{\beta_1}/\varphi(M)$. Their
sum vanishes unless $\chi_1(u)=\left(\frac us\right)$ for every $u\equiv1\pmod{m_0}$ prime to $s$, and
then $\beta_1\leq1-c_0/\log Z$ because $s$ is not bad, so that
$|\operatorname{li}U^{\beta_1}|\leq r\operatorname{li}U$ with $r=4\log U\,e^{-c_0\log U/\log Z}$.
Since $\sum_b|m_M(b,U)|\leq2(|P|+U^{1/2})$ when $k\leq1/2$, and $\sum_b\pi(\sqrt U;M,b)\leq U^{1/2}$,
this gives
\[
 \Bigl|\sum_{q\in P}\Bigl(\frac qs\Bigr)\Bigr|\leq(4r+2k)(|P|+U^{1/2})+U^{1/2}.
\]

It remains to bound the right side. Theorem~\ref{hf-primes-progressions}(a) for $M=m_0$, with the
bound for $\beta_1$ in part (c), shows that $|P|\geq U^{9/10}$ for $Y$ large, so
$U^{1/2}\leq U^{-2/5}|P|$. Since $\log Z\leq4\log U/\log^2Y+A\log Y$ and
$Y^{\varepsilon/6}\leq\log U\leq Y^{\varepsilon/5}\log^2Y$, for $Y$ large the quantities $4r$, $2k$
and $2U^{-2/5}$ are each at most $Y^{-B}/8$. The bound follows.
\end{proof}

\begin{lemma}\label{hf-prime-power-sums}
For real $y\geq3$ put $c_y(n)=\Lambda(n)\min(1,y/n)$. For integers $N,m_0\geq1$ and a prime $p$ let
$s_N(p)$ be the sum of $c_y(n)$ over the powers $n\leq N$ of $p$, and let $R_N(m_0)$ be the sum of
$c_y(n)$ over the prime powers $n=p^k\leq N$ with $k$ even or $p\mid m_0$. There is $C>0$ such that
for all $y\geq3$ and $N,m_0\geq1$:
\begin{enumerate}
\item $s_N(p)\leq3\log y$ for every prime $p$;
\item $\sum_{p\leq N}s_N(p)^2\leq Cy\log y$;
\item $R_N(m_0)\leq C(\sqrt y+w\log y)$ for every $w\geq1$ such that every prime $p\mid m_0$ with
  $p\leq N$ satisfies $p\leq w$;
\item $\sum_{n\leq Z}c_y(n)\leq Cy\log Z$ for every real $Z\geq2$.
\end{enumerate}
\end{lemma}
\leanref{Development/AutomorphicArithmetic/MeanSquareBounds/SumMinLeDiv.lean}{QH.exists_primePowerSum_le_mul_log}
\leanref{Development/AutomorphicArithmetic/MeanSquareBounds/SPsqLe.lean}{QH.sPsq_le}
\leanref{Development/AutomorphicArithmetic/MeanSquareBounds/RRLe.lean}{QH.RR_le}
\leanref{Development/AutomorphicArithmetic/MeanSquareBounds/SumCNLe.lean}{QH.sum_cN_le}

\begin{proof}
For (a), if $p\leq y$ the powers $p^k\leq y$ contribute at most $\log y$ and the others at most
$\sum_{p^k>y}y\log p/p^k\leq2\log p$; if $p>y$, then $s_N(p)\leq2y\log p/p\leq2\log y$, since
$t\mapsto\log t/t$ decreases for $t\geq3$. For (b), the primes $p\leq y$ contribute at most
$3\log y\sum_{p\leq y}s_N(p)\ll y\log y$ by (a) and Chebyshev's bound, and the primes $p>y$
contribute at most $\sum_{p>y}4y^2\log^2p/p^2\ll y\log y$. For (c), the prime powers with even
exponent contribute $\ll\sum_{p^{2j}\leq y}\log p+\sum_{p^{2j}>y}y\log p/p^{2j}\ll\sqrt y$, and the
at most $w$ primes $p\mid m_0$ with $p\leq N$ contribute at most $3w\log y$ by (a). For (d),
$\sum_{n\leq Z}c_y(n)\leq y\sum_{n\leq Z}\Lambda(n)/n\ll y\log Z$ by Mertens' estimate.
\end{proof}

\begin{proposition}\label{hf-mean-square}
For $Z,y>0$, a prime $q$ and $\kappa\colon\mathbf N\to\mathbf C$ put
\[
 S(Z,y,q,\kappa)=\sum_{n\leq Z}\Lambda(n)\min(1,y/n)\,\kappa(n)\left(\tfrac nq\right).
\]
For integers $m_0,a_0$ and $U>0$ let $P_{m_0,a_0}(U)$ be the set of primes $q$ with $\sqrt U<q\leq U$
and $q\equiv a_0\pmod{m_0}$. There is $C>0$ such that for all $0<\varepsilon\leq1$ and $K\geq1$
there is $Y_0$ with the following property. Let $Y\geq Y_0$, let $m_0\geq1$ and $a_0$ be integers
with $8\mid m_0$ and $\gcd(a_0,m_0)=1$, and let $U$ satisfy $\tfrac14\log^2Y\log m_0\leq\log U$ and
$Y^{\varepsilon/6}\leq\log U\leq Y^{\varepsilon/5}\log^2Y$. Suppose that every prime divisor $p$ of
$m_0$ has $p\leq Y^{\varepsilon/5}$ or $p>Y^{400K}$. Then for every $y\in[Y,Y^K]$ and every $\kappa$
with $|\kappa(n)|\leq1$ for all $n$,
\[
 \sum_{q\in P_{m_0,a_0}(U)}|S(Y^{400K},y,q,\kappa)|^2\leq C\,y\log y\,|P_{m_0,a_0}(U)|.
\]
\end{proposition}
\leanref{Development/AutomorphicArithmetic/MeanSquareBounds/meanSqK.lean}{QH.meanSqK}

\begin{proof}
Let $C_1$ be the constant of Lemma~\ref{hf-prime-power-sums}. Let $Y$ be large, write
$P=P_{m_0,a_0}(U)$, $Z=Y^{400K}$ and $N_0=\lfloor Z\rfloor$, and fix $y\in[Y,Y^K]$ and $\kappa$.
Then $S(Z,y,q,\kappa)=\sum_{n\leq N_0}c_y(n)\kappa(n)\left(\frac nq\right)$ with $c_y$ as in
Lemma~\ref{hf-prime-power-sums}.

\emph{Step 1.} Since $\log U\geq Y^{\varepsilon/6}\geq1600K\log Y$ for $Y$ large, every $q\in P$
satisfies $q>\sqrt U\geq Y^{800K}=Z^2\geq N_0^2$.

\emph{Step 2.} Expanding the square and using $|\kappa(n)|\leq1$,
\[
 \sum_{q\in P}|S(Z,y,q,\kappa)|^2\leq\sum_{n,n'\leq N_0}c_y(n)c_y(n')\Bigl|\sum_{q\in P}
 \Bigl(\frac nq\Bigr)\Bigl(\frac{n'}q\Bigr)\Bigr|.
\]

\emph{Step 3.} For an integer $N'\geq1$ let $s(N')$ be the product of the primes $p\nmid m_0$ that
divide $N'$ to an odd power. Let $q$ be a prime with $q\equiv a_0\pmod{m_0}$ and $q>N'$. Write
$N'=2^\alpha m'f^2s(N')$ with every prime factor of $m'$ dividing $m_0$. Since $8\mid m_0$, the
symbols $\left(\frac2q\right)$ and $\left(\frac{p}{q}\right)$ for odd $p\mid m_0$ depend only on $q$
modulo $m_0$, and by quadratic reciprocity
$\left(\frac{s(N')}q\right)=\pm\left(\frac q{s(N')}\right)$ with a sign depending only on $q$ modulo
$4$. Hence $\left(\frac{N'}q\right)=\sigma\left(\frac q{s(N')}\right)$ with $\sigma=\pm1$ depending
only on $N'$, $m_0$ and $a_0$. With $N'=nn'<q$ (Step 1) this gives
\[
 \Bigl|\sum_{q\in P}\Bigl(\frac nq\Bigr)\Bigl(\frac{n'}q\Bigr)\Bigr|
 =\Bigl|\sum_{q\in P}\Bigl(\frac q{s(nn')}\Bigr)\Bigr|.
\]

\emph{Step 4.} The integer $s(nn')$ is odd, squarefree, prime to $m_0$ and at most
$N_0^2\leq Y^{800K}$. Let $s^*$ be the integer of Lemma~\ref{hf-twisted-prime-sums} with $A=800K$
and $B=K+1$. Then the inner sum in Step 3 is at most $|P|$ in absolute value, and at most
$Y^{-K-1}|P|$ when $s(nn')\notin\{1,s^*\}$. So
\[
 \sum_{q\in P}|S(Z,y,q,\kappa)|^2\leq\Bigl(D_1+D_{s^*}+Y^{-K-1}\Bigl(\sum_{n\leq N_0}c_y(n)\Bigr)^2\Bigr)|P|,
 \qquad D_r=\sum_{\substack{n,n'\leq N_0\\ s(nn')=r}}c_y(n)c_y(n').
\]

\emph{Step 5.} The function $c_y$ is supported on prime powers. Let $n=p^a$ and $n'=p'^b$ be prime
powers with $c_y(n)c_y(n')\neq0$. If $s(nn')=1$, then either $p=p'$, or each of $n,n'$ has even
exponent or a prime dividing $m_0$; hence $D_1\leq\sum_{p\leq N_0}s_{N_0}(p)^2+R_{N_0}(m_0)^2$. Let
$r>1$. If $r$ is prime and $s(nn')=r$, then both $n,n'$ are powers of $r$, or one is a power of $r$
and the other is counted in $R_{N_0}(m_0)$. If $r$ is not prime and $s(nn')=r$, then $r=p_1p_2$ with
$n$ a power of $p_1$ and $n'$ a power of $p_2$ or the reverse. With $B'=\max_ps_{N_0}(p)$ this gives
$D_r\leq2B'^2+2B'R_{N_0}(m_0)$.

\emph{Step 6.} By Lemma~\ref{hf-prime-power-sums}: $B'\leq3\log y$;
$\sum_ps_{N_0}(p)^2\leq C_1y\log y$; and, with $w=Y^{\varepsilon/5}$,
$R_{N_0}(m_0)\leq C_1(\sqrt y+w\log y)\leq5C_1\sqrt y$. Here part (c) applies because a prime
$p\mid m_0$ with $p\leq N_0\leq Y^{400K}$ has $p\leq Y^{\varepsilon/5}$ by hypothesis, and
$w\log y\leq y^{1/5}\log y\leq4\sqrt y$. Finally
$\sum_{n\leq N_0}c_y(n)\leq C_1y\log Z=400C_1Ky\log Y$, and $(400K\log Y)^2\leq Y$ for $Y$ large, so
$Y^{-K-1}(\sum_nc_y(n))^2\leq C_1^2y^2Y^{-K}\leq C_1^2y$.

\emph{Step 7.} Using $\log y\geq1$, $\sqrt y\log y\leq2y$ and $\log^2y\leq2y\log y$, Steps 4 to 6
give $\sum_{q\in P}|S(Z,y,q,\kappa)|^2\leq C\,y\log y\,|P|$ with
$C=C_1+25C_1^2+36+60C_1+C_1^2$.
\end{proof}

\begin{lemma}\label{hf-log-derivative-box}
For every $\delta>0$ there is $K_\delta\geq0$ with the following property. Let $Q\geq2$, let $\chi$
be a nontrivial Dirichlet character of modulus $N\leq Q$, and suppose that $L(s,\chi)\neq0$ for
$\operatorname{Re}s>1-1/400$ and $|\operatorname{Im}s|\leq Q$. Then
$\left|\frac{L'}{L}(s,\chi)\right|\leq K_\delta\log Q$ for every $s$ with
$\operatorname{Re}s\geq1-1/400+\delta$ and $|\operatorname{Im}s|\leq Q-1$.
\end{lemma}
\leanref{Formalism/GL3MeanSquares/GenusMeanAsymptotics/JacobiCharacterSums/TailBox.lean}{Formalism.HermitianLattices.Ext.BoxTail.tail_box_negLD}

\begin{proof}
There is $C_0$ with $\left|\frac{L'}{L}(s,\chi)\right|\leq\sum_n\Lambda(n)n^{-\sigma}\leq
1/(\sigma-1)+C_0$ for $\sigma=\operatorname{Re}s>1$; this gives the bound for $\sigma\geq11/10$. For
$\sigma<11/10$ apply Landau's lemma \cite[Section 6.1]{MV07} to $L(\cdot,\chi)$ on a disc of bounded
radius about $11/10+i\operatorname{Im}s$, using the polynomial bound
$|L(w,\chi)|\leq(N(|\operatorname{Im}w|+2))^{O(1)}$ in the strip and the lower bound for
$|L(w,\chi)|$ at $\operatorname{Re}w=11/10$. It writes $\frac{L'}{L}(s,\chi)$ as the sum of
$1/(s-\rho)$ over $O(\log Q)$ zeros $\rho$ near $s$, plus $O(\log Q)$. These zeros have
$|\operatorname{Im}\rho|\leq Q$, so $\operatorname{Re}\rho\leq1-1/400$ and $|s-\rho|\geq\delta$.
Hence $\left|\frac{L'}{L}(s,\chi)\right|\ll(1+1/\delta)\log Q$.
\end{proof}

\begin{lemma}\label{hf-box-tail}
For a Dirichlet character $\chi$ and $y\geq1$ put
$\Sigma_\chi(y)=\sum_{n\leq y}\Lambda(n)\chi(n)(y/n-1)$. For every $0<\delta\leq1/800$ there is
$A_\delta\geq0$ with the following property. Let $Q\geq2$, let $\chi$ be a nontrivial Dirichlet character
of modulus $N\leq Q$ such that $L(s,\chi)\neq0$ for $\operatorname{Re}s>1-1/400$ and
$|\operatorname{Im}s|\leq Q$, and let $z\geq1$ and $0<\eta\leq1$. Then
\[
 \Bigl|\Sigma_\chi(z+z\eta)-\Sigma_\chi(z)+z\eta\frac{L'}{L}(1,\chi)\Bigr|\leq
 A_\delta\log Q\bigl(\eta z^{1-1/400+\delta}\log Q+z^2/Q\bigr).
\]
\end{lemma}
\leanref{Formalism/GL3MeanSquares/GenusMeanAsymptotics/JacobiCharacterSums/TailContour.lean}{Formalism.HermitianLattices.Ext.BoxTail.tail_contour}

\begin{proof}
Let $C_0$ be a constant with $\left|\frac{L'}{L}(w,\chi)\right|\leq1/(\operatorname{Re}w-1)+C_0$
for $\operatorname{Re}w>1$. Put $\sigma_1=1-1/400+\delta$, $T=Q/2$ and
$F_y(w)=-\frac{L'}{L}(w,\chi)\,y^w/(w(w-1))$. By
Perron's formula, $\Sigma_\chi(y)=\frac1{2\pi i}\int_{(2)}F_y(w)\,dw$ for $y\geq1$. Move the part of
the line of integration with $|\operatorname{Im}w|\leq T$ to $\operatorname{Re}w=\sigma_1$. The
rectangle contains no zero of $L(w,\chi)$, and the only pole of $F_y$ in it is at $w=1$, with residue
$-y\frac{L'}{L}(1,\chi)$. So $2\pi\bigl(\Sigma_\chi(y)+y\frac{L'}{L}(1,\chi)\bigr)$ is the sum of
the integral of $F_y$ over $\operatorname{Re}w=2$, $|\operatorname{Im}w|>T$, of the two horizontal
segments $\sigma_1\leq\operatorname{Re}w\leq2$, $\operatorname{Im}w=\pm T$, and of the segment
$\operatorname{Re}w=\sigma_1$, $|\operatorname{Im}w|\leq T$. Subtract these identities for
$y=z+z\eta$ and $y=z$; both values of $y$ are at most $2z$.

On $\operatorname{Re}w=2$, $|F_y(w)|\leq(1+C_0)y^2/t^2$ with $t=\operatorname{Im}w$, so the two
vertical tails contribute $O((1+C_0)z^2/Q)$. On the horizontal segments
$|F_y(w)|\leq K_\delta\log Q\cdot y^2/T^2$ by Lemma~\ref{hf-log-derivative-box}, so they contribute
$O(K_\delta z^2\log Q/Q)$. On the line $\operatorname{Re}w=\sigma_1$,
$|(z+z\eta)^w-z^w|\leq2\eta|w|z^{\sigma_1}$ and $|w-1|\geq1/800$, so the difference of the two
integrals is at most
$K_\delta\log Q\cdot4\eta z^{\sigma_1}\int_{-T}^T|w-1|^{-1}\,dt\ll K_\delta\eta z^{\sigma_1}\log^2Q$.
Adding the three bounds gives the claim with $A_\delta=160K_\delta+32(1+C_0)$.
\end{proof}

\begin{lemma}\label{hf-chebyshev-from-twisted-sum}
There is $C>0$ such that for every $0<\varepsilon\leq1/4$, $\tau>0$ and $K\geq1$ there is $Y_0$ with
the following property. Let $Y\geq Y_0$ and $Q$ with $Y^{\varepsilon/6}\leq\log Q\leq
Y^{\varepsilon/4}$, let $q$ be a prime, let $\kappa\colon\mathbf N\to\mathbf C$, and let $\chi$ be a
nontrivial Dirichlet character of modulus $N\leq Q$ with $\chi(n)=\kappa(n)\left(\frac nq\right)$
for all $n$ and $L(s,\chi)\neq0$ for $\operatorname{Re}s>1-1/400$, $|\operatorname{Im}s|\leq Q$.
If $|S(Y^{400K},y,q,\kappa)|\leq y^{1/2+\tau}/2$ for every $y\in[Y,Y^K]$, then
$\left|\frac{L'}{L}(1,\chi)\right|\leq C\log Y$ and $|g_\chi(y)|\leq y^{1/2+\tau}$ for every
$y\in[Y,Y^K]$.
\end{lemma}
\leanref{Development/AutomorphicArithmetic/MeanSquareBounds/BoxOfS.lean}{QH.box_of_S}

\begin{proof}
Let $C_0$ be the constant of Mertens' bound $\sum_{n\leq y}\Lambda(n)/n\leq C_0\log y$ for $y\geq2$,
put $C=400C_0+1$, and let $Y$ be large. For $z\geq1$ put
$T_\chi(z)=-\frac{L'}{L}(1,\chi)-\sum_{n\leq z}\Lambda(n)\chi(n)/n$.

\emph{Step 1: the tail.} Let $X=Y^K$, $z=X^{400}$ and $h=z/X^2$, so $h=z\eta$ with
$\eta=X^{-2}\leq1$. Then
\[
 hT_\chi(z)=\Bigl(\Sigma_\chi(z+h)-\Sigma_\chi(z)-h\sum_{n\leq z}\frac{\Lambda(n)\chi(n)}n\Bigr)
 -\Bigl(\Sigma_\chi(z+h)-\Sigma_\chi(z)+h\frac{L'}{L}(1,\chi)\Bigr),
\]
with $\Sigma_\chi$ as in Lemma~\ref{hf-box-tail}. The first bracket is
$\sum_{z<n\leq z+h}\Lambda(n)\chi(n)((z+h)/n-1)$, which has at most $h+1$ terms, each at most
$h\log(z+h)/z$. The second bracket is at most
$A_\delta\log Q(\eta z^{1-1/400+\delta}\log Q+z^2/Q)$ by Lemma~\ref{hf-box-tail} with
$\delta=\varepsilon/(1600K)$, since $Q\geq e$. Dividing by $h$ and multiplying by $X$,
\[
 X|T_\chi(z)|\leq A_\delta\log^2Q\cdot Y^{\varepsilon/4}+A_\delta\log Q\cdot X^{403}/Q+(X^{398}+1)\log(z+h)/X^{399}.
\]
Since $\log Q\leq Y^{\varepsilon/4}$, the first term is at most $A_\delta Y^{3\varepsilon/4}\leq Y^\varepsilon/3$
for $Y$ large. Since $\log Q\geq Y^{\varepsilon/6}$ and $\log Q\leq X$, the second is at most
$A_\delta Y^{404K}\exp(-Y^{\varepsilon/6})\leq1/3$ for $Y$ large. The third is at most $802$ for $X\geq3$. So
$X|T_\chi(z)|\leq Y^\varepsilon$ once $Y^\varepsilon\geq3\cdot802$.

\emph{Step 2: the Chebyshev sum.} For $\chi\neq1$ and $y\leq z$, splitting each term of $g_\chi(y)$
according to $n\leq y$ gives
\[
 g_\chi(y)=-\sum_{n\leq z}\Lambda(n)\chi(n)\min(1,y/n)-yT_\chi(z).
\]
Since $\chi(n)=\kappa(n)\left(\frac nq\right)$, the sum is $S(Y^{400K},y,q,\kappa)$. For
$y\in[Y,Y^K]$, Step 1 gives $|yT_\chi(z)|\leq X|T_\chi(z)|\leq Y^\varepsilon$, and
$Y^\varepsilon\leq Y^{1/4}\leq y^{1/2}/2\leq y^{1/2+\tau}/2$ because $\varepsilon\leq1/4$ and
$Y\geq16$. With the hypothesis on $S$ this gives $|g_\chi(y)|\leq y^{1/2+\tau}$.

\emph{Step 3: the logarithmic derivative.} Step 1 with $K=1$ gives
$Y|T_\chi(Y^{400})|\leq Y^\varepsilon\leq Y$, so $|T_\chi(Y^{400})|\leq1$. Hence
\[
 \Bigl|\frac{L'}{L}(1,\chi)\Bigr|\leq|T_\chi(Y^{400})|+\sum_{n\leq Y^{400}}\frac{\Lambda(n)}n
 \leq1+400C_0\log Y\leq C\log Y.
\]
\end{proof}

\begin{lemma}\label{hf-most-primes-good}
There is $C>0$ such that for all $0<\varepsilon<1/52$, $\tau>0$, $K\geq1$, $0<\eta$ and integers
$N_0\geq1$ there is $Y_0$ with the following property. Let $Y\geq Y_0$, put $z=Y^{\varepsilon/5}$,
and let $z'$ satisfy $z\log^2Y\leq z'\leq K^2z\log^2Y$. Let $P_1$ be the set of primes $q$ with
$e^{z/2}<q\leq e^z$ and $q\equiv3\pmod8$, and for $q\in P_1$ let $P_2(q)$ be the set of primes $d$
with $e^{z'/2}<d\leq e^{z'}$, $d\equiv7\pmod{16}$, $q\mid d+1$ and $p\mid d+1$ for every odd prime
$p\leq z$. Call a primitive character $\chi$ \emph{good at level} $Q$ if $\chi\neq1$, its conductor is
at most $Q$, $L(s,\chi)\neq0$ for $\operatorname{Re}s>1-1/400$ and $|\operatorname{Im}s|\leq Q$,
$\left|\frac{L'}{L}(1,\chi)\right|\leq C\log Y$ and $|g_\chi(y)|\leq y^{1/2+\tau}$ for all
$y\in[Y,Y^K]$. Then $P_1$ is nonempty, and there is $G_1\subset P_1$ with $|G_1|\geq(1-\eta)|P_1|$
such that for every $q\in G_1$:
\begin{enumerate}
\item for every character $\xi$ modulo $N_0$, every primitive character inducing
  $a\mapsto\xi(a)\left(\frac aq\right)$ on the integers prime to $N_0q$ is good at level $N_0e^{z}$;
\item $P_2(q)$ is nonempty, and there is $G_2\subset P_2(q)$ with $|G_2|\geq(1-\eta)|P_2(q)|$ such
  that for every $d\in G_2$, every character $\xi$ modulo $N_0$ and every $i\in\{0,1\}$, every
  primitive character inducing $a\mapsto\xi(a)\left(\frac aq\right)^i\left(\frac ad\right)$ on the
  integers prime to $N_0qd$ is good at level $N_0e^{z+z'}$.
\end{enumerate}
\end{lemma}
\leanref{Development/AutomorphicArithmetic/MeanSquareBounds/PairShare.lean}{QH.pair_share}

\begin{proof}
Both parts are instances of one statement about primes $q$ in a progression $P_{m_0,a_0}(U)$ with
$\log U$ between $Y^{\varepsilon/6}$ and $Y^{\varepsilon/5}\log^2Y$, applied with
$(m_0,a_0,U)=(8,3,e^z)$ for $q$ and, for $d$, with $m_0$ equal to $16$ times $q$ times the odd
primes up to $z$, $a_0\equiv7\pmod{16}$ and $a_0\equiv-1$ modulo the other factors, and $U=e^{z'}$.
The number of primes in $P_{m_0,a_0}(U)$ is at least a constant times $U/(\varphi(m_0)\log U)$ by
the prime number theorem in arithmetic progressions with Siegel's bound, so the sets are nonempty
and large.

A prime $q$ fails (a) for one of three reasons. First, some character inducing
$\xi\left(\frac{\cdot}{q}\right)$ has a zero in the box of height $Q=N_0U$; since such a character has
conductor $fq$ with $f\leq N_0$, distinct $q$ give distinct characters, and
Theorem~\ref{hf-zero-density-box} bounds the number of such $q$ by $CQ^{1/4}$, which is at most
$\frac\eta3|P_{m_0,a_0}(U)|$ for $Y$ large. Second, $|S(Y^{400K},y,q,\xi)|>y^{1/2+\tau}/2$ for some
$y\in[Y,Y^K]$; by Proposition~\ref{hf-mean-square} at a net of values of $y$ and Chebyshev's
inequality there are at most $\frac\eta3|P_{m_0,a_0}(U)|$ such $q$ for each of the boundedly many
$\xi$. Third, neither happens, and then Lemma~\ref{hf-chebyshev-from-twisted-sum} shows that the
character is good. The same argument applies to $d$.
\end{proof}

\begin{theorem}\label{hf-good-pairs-exist}
Let $M\geq1$ be an integer and let $0<\varepsilon<1/52$ and $A\geq1$. There is $Y_1$ such that for
every $Y\geq Y_1$ there are primes $q,d$ such that $(q,d)$ is good for $(M,\varepsilon,Y,Y^A)$
(Definition~\ref{hf-good-pair}).
\end{theorem}
\leanref{Development/AutomorphicArithmetic/MeanSquareBounds/pairEx.lean}{QH.pairEx}
\leanref{Development/AutomorphicArithmetic/CharacterBoxes/PairGood.lean}{UpperTheory.Top.PairEx}

\begin{proof}
Apply Lemma~\ref{hf-most-primes-good} with $\tau=\varepsilon/2$, $K=A$, $\eta=1/2$, $N_0=8M$ and
$z'=Y^{\varepsilon/5}\log^2Y$, and pick $q\in G_1$ and $d\in G_2$. With $z=Y^{\varepsilon/5}$, the
conditions defining $P_1$ and $P_2(q)$ say that $(q,d)$ is admissible at $Y$. For $Y$ large,
$q>e^{z/2}>8M$ and $d>q$, so $\gcd(M,qd)=1$.

Let $\xi$ be a character modulo $M$ and $(i,j)\neq(0,0)$. Lift $\xi$ to a character modulo $8M$.
If $j=0$, then $\xi\left(\frac{\cdot}{q}\right)$ is a character modulo $8Mq$, so the primitive
character $\chi$ inducing it on the integers prime to $2Mqd$ also induces it on the integers prime to
$8Mq$, and (a) of Lemma~\ref{hf-most-primes-good} applies. If $j=1$, part (b) applies. In both cases
$|g_\chi(y)|\leq y^{1/2+\varepsilon/2}$ for $y\in[Y,Y^A]$, and $\chi$ has no zeros in the box of
height $8Me^{z}e^{jz'}$; this gives (c) of Definition~\ref{hf-good-pair} when $\xi$ is a real
primitive character of conductor dividing $8M$. The logarithm of that height $Q$ lies between
$Y^{\varepsilon/6}$ and $3Y^{\varepsilon/5}\log^2Y$. A nontrivial character of modulus at most $Q$
whose $L$-function has no zeros in the box of height $Q$ satisfies
$\left|\frac{L'}{L}(1,\chi)\right|\leq1000\log\log Q+O(1)$, which is at most $4\log Y\leq4\log y$
for $Y$ large since $\varepsilon<1/52$.
\end{proof}

\begin{lemma}\label{hf-bounded-characters}
Let $\psi$ be a primitive Dirichlet character of conductor $N_\psi$, let $\Phi$ be a Galois datum, and let
$(q,d,Q_S,E,\gamma,j,\operatorname{art})$ be a host field on $\Phi$ with $N=Q_Sqd$. Suppose that
$N_\psi\mid Q_S$, that every prime divisor of $Q_S$ divides $2N_\psi$, and that the host field has
property (e) of Proposition~\ref{hf-host-field-exists}. Let $\varepsilon>0$ and $1<Y\leq x$, and
suppose that $(q,d)$ is good for $(N_\psi,\varepsilon,Y,x)$. Let $\theta$ be a character of $G_\Phi$
that lies in $X^-_\Phi$, or lies in $X^+_\Phi$ with $\chi_\theta$ inducing no power of $\psi$. Then
$|g_{\chi_\theta}(x)|\leq x^{1/2+\varepsilon}$ and $\left|\frac{L'}{L}(1,\chi_\theta)\right|\leq
4\log x$.
\end{lemma}
\leanref{Development/AutomorphicArithmetic/CharacterBoxes/BoundedTw.lean}{UpperTheory.Top.boundedTw}

\begin{proof}
By property (e) there are $\xi\in\langle\psi\rangle$ and $i,j\in\{0,1\}$ with
$\theta(\iota_d)=(-1)^j$ and $\theta(\operatorname{art}(a))=\xi(a)\left(\frac aq\right)^i
\left(\frac ad\right)^j$ for $a$ prime to $N$. If $i=j=0$, then $\theta(\iota_d)=1$, so
$\theta\in X^+_\Phi$, and $\chi_\theta$ induces $\xi$, which is a power of $\psi$; this is excluded.
So $(i,j)\neq(0,0)$. Every prime dividing $N$ divides $2N_\psi qd$, so $\chi_\theta$ induces
$\xi(a)\left(\frac aq\right)^i\left(\frac ad\right)^j$ on the integers prime to $2N_\psi qd$.
Condition (b) of Definition~\ref{hf-good-pair} at $y=x$ gives the two bounds.
\end{proof}

\section{Windows of scales}\label{section-hf-windows}

\begin{definition}\label{hf-windows}
A \emph{scale} is a triple $s=(x,\varepsilon,J)$ of real numbers with $x>1$, $0<\varepsilon<1/52$
and $J\geq1$. With it we use $\delta_1=x^{-\varepsilon/8}$, $b=\exp(x^{\varepsilon})$ and
$X'=\exp(x^{\varepsilon/3})$.

Let $\lambda\in\mathbf R$. A \emph{set of scales with windows of exponent} $\lambda$ is a pair
$(J_W,W)$ of a real number $J_W\geq1$ and a set $W\subset\mathbf R$ with windows of exponent
$\lambda$ (Definition~\ref{wl-windows}). A scale $(x,\varepsilon,J)$ \emph{belongs} to it if $J=J_W$ and
$x\in W$.
\end{definition}
\leanref{Formalism/LocalRingCardinality/DataParams₂.lean}{Formalism.BasicNotions.Scale}
\leanref{Development/AutomorphicArithmetic/PrimeDiscriminants/Scales.lean}{GlobalTheory.EX.Scales}

\begin{proposition}\label{hf-windows-exist}
Let $\psi$ be a primitive Dirichlet character of conductor $N_\psi$, let $0<\varepsilon<1/52$ and
$c\geq1$, and suppose that $\psi$ satisfies the weakening of Vandiver's conjecture with exponent
$401$ (Definition~\ref{sha-vandiver-weakening}). Then there is $J\geq1$ such that for every $y$ there is $Z\geq y$ with the following property. For
every $x\in[Z,Z^{3/2}]$ there are a real number $Y$, primes $q,d$ and an integer $e\geq0$ such that:
\begin{enumerate}
\item $x^{1/J}\leq Y\leq x$, and $(q,d)$ is good for $(N_\psi,\varepsilon,Y,x)$
  (Definition~\ref{hf-good-pair});
\item $x^{2+8\varepsilon}\leq c\,2^e\leq2x^{2+8\varepsilon}$;
\item every host field $E$ of $(\psi,e,q,d)$ (Definition~\ref{hf-host-field-of-character})
  satisfies the Tate-Shafarevich bound $\bigl(\lfloor x^{400}\rfloor,\exp(x^{\varepsilon}/2)\bigr)$
  (Definition~\ref{sha-bounded-field}).
\end{enumerate}
\end{proposition}
\leanref{Development/AutomorphicArithmetic/CharacterBoxes/WinPairs.lean}{UpperTheory.Top.winPairs}
\leanref{Development/AutomorphicArithmetic/CharacterBoxes/TwistForm.lean}{UpperTheory.Top.WinDatum}
\leanref{Development/AutomorphicArithmetic/CharacterBoxes/ExistsPageN.lean}{UpperTheory.Top.VWinAt}

\begin{proof}
For $x\geq1$ let $e(x)$ be an integer with $x^{2+8\varepsilon}\leq2^{e(x)}\leq2x^{2+8\varepsilon}$.
The hypothesis, applied to $\varepsilon$ and $e$, gives $J\geq1$ and $X_0$. Theorem~\ref{hf-good-pairs-exist}
with $M=N_\psi$ and $A=J$ gives $Y_1$. Put $k=c^{1/(2+8\varepsilon)}\geq1$.

Let $y$ be given. Choose $M'\geq\max(1,y,Y_1,k^3,k^{400})$ and $X\geq\max(X_0,M'^J)$. Then
$X^{1/J}\geq M'\geq Y_1$, so there is a pair $(q,d)$ good for $(N_\psi,\varepsilon,X^{1/J},X)$; it is
admissible at $X^{1/J}$. The hypothesis gives $W$ with $X^{1/J}\leq W$ and $W^2\leq X$ such that for
every $x'\in[W,W^2]$, every host field of $(\psi,e(x'),q,d)$ satisfies the Tate-Shafarevich bound
$\bigl(\lfloor x'^{401}\rfloor,\exp(x'^{\varepsilon}/2)\bigr)$. Put $Z=kW$, so $Z\geq W\geq M'\geq y$.

Let $x\in[Z,Z^{3/2}]$ and write $x=kx'$. Then $W\leq x'\leq x$, and $x\leq(kW)^{3/2}\leq W^2\leq X$
because $k^3\leq W$. Take $Y=X^{1/J}$ and $e=e(x')$. Then $x^{1/J}\leq Y\leq W\leq x$, and the pair
stays good when the interval $[Y,X]$ is shortened to $[Y,x]$. Since
$(kx')^{2+8\varepsilon}=c\,x'^{2+8\varepsilon}$, the choice of $e(x')$ gives (b). Finally
$x'\in[W,W^2]$, $x^{400}=k^{400}x'^{400}\leq x'^{401}$ because $k^{400}\leq W\leq x'$, and
$\exp(x'^\varepsilon/2)\leq\exp(x^\varepsilon/2)$. The Tate-Shafarevich bound $(t,B)$ implies the
bound $(t',B')$ for $t'\leq t$ and $B'\geq B$, which gives (c).
\end{proof}

\section{The Behrend set}\label{section-hf-behrend}

\begin{definition}\label{hf-no-short-relations}
Let $T$ be a finite set of non-negative integers and let $L\geq1$ be an integer. The set $T$ has
\emph{no relations of length} $L$ if there are no pairwise distinct $s_1,s_2,s_3\in T$ and integers
$i\neq j$ with $0<|i|,|j|\leq L$ and $j(s_2-s_1)=i(s_3-s_1)$.
\end{definition}
\leanref{Formalism/FrobeniusOrbits/HostThickExists.lean}{Formalism.FrobeniusOrbits.NoRel}

\begin{theorem}\label{hf-behrend-set}
There is $C>0$ such that for all integers $n\geq0$ and $L\geq2$ there is a finite set $T$ of
non-negative integers with no relations of length $L$ (Definition~\ref{hf-no-short-relations}),
with $4Ls<n$ for every $s\in T$, and with
\[
 |T|\geq\frac nL\exp\bigl(-C\sqrt{\log n\log L}\bigr).
\]
\end{theorem}
\leanref{Formalism/FrobeniusOrbits/FrobeniusOrbitsProofs/ThickExists.lean}{Formalism.FrobeniusOrbits.hostThickExists}
\leanref{Formalism/FrobeniusOrbits/HostThickExists.lean}{Formalism.FrobeniusOrbits.HostThickExists}

This is Behrend's construction \cite{Beh46}, with $C=32$. The integers are the points of a sphere in
$\{0,\ldots,D-1\}^S$, read as digits in base $4LD$. A relation $j(s_2-s_1)=i(s_3-s_1)$ with
$|i|,|j|\leq L$ involves no carries, so it holds digit by digit and puts three distinct points of
the sphere on a line, which is impossible.

\begin{corollary}\label{hf-behrend-set-at-scale}
There is $C_T>0$ such that for all $0<\varepsilon<1/52$ and $C_o,C_L$ there is $x_0$ with the
following property. Let $x\geq x_0$, let $\Phi$ be a Galois datum (Definition~\ref{hf-frame}) whose group
has order $n$ with $x^{2+8\varepsilon}/2\leq n\leq C_ox^{2+8\varepsilon}$ and whose index $o$ is at
most $C_o$, and let $L$ be an integer with $2\leq L\leq C_L\log x$. Then there is a finite set $T$ of
non-negative integers with no relations of length $L$ such that $4Ls<n_0$ for $s\in T$, $4L|T|\leq
n_0$, $|T|\geq x^{-\varepsilon}n$ and $|T|\geq\frac{n_0}{L}\exp(-C_T\sqrt{\log n_0\log L})$.
\end{corollary}
\leanref{Development/AutomorphicArithmetic/ThickFrames/ThickFrame2Large.lean}{UpperTheory.BernoulliSumsFrobeniusOrbits.thickFrame2}

\begin{proof}
Let $C$ be the constant of Theorem~\ref{hf-behrend-set} and put $C_T=2C$. Let $x$ be large and
$n_0=n/o$, so that $x^{2+8\varepsilon}/(2C_o)\leq n_0\leq n$.

\emph{Step 1: the set.} Theorem~\ref{hf-behrend-set} for $n_0$ and $L$ gives a set $T'$ with no
relations of length $L$, with $4Ls<n_0$ for $s\in T'$ and with
$|T'|\geq(n_0/L)e^{-Cv}$, where $v=\sqrt{\log n_0\log L}$. Let $T=T'\setminus\{0\}$. It has no
relations of length $L$ and $4Ls<n_0$ for $s\in T$. Its elements are distinct integers $s$ with
$1\leq s<n_0/(4L)$, so $|T|\leq\max T<n_0/(4L)$ and $4L|T|\leq n_0$.

\emph{Step 2: the Behrend bound.} We have $v\geq\log2$, since $\log L\geq\log2$ and
$\log n_0\geq\log2$ for $x$ large. Then $(n_0/L)(e^{-Cv}-e^{-2Cv})\geq(n_0/L)e^{-2Cv}(e^{C\log2}-1)$,
and $(n_0/L)e^{-2Cv}\to\infty$ as $x\to\infty$ because $v\ll\sqrt{\log x\log\log x}$. So
$|T|\geq|T'|-1\geq(n_0/L)e^{-C_Tv}$ for $x$ large.

\emph{Step 3: the bound in terms of $x$.} Put $\delta=\varepsilon/(3C)$. By the inequality of
arithmetic and geometric means, $v\leq(\delta\log n_0+\log L/\delta)/2$, so
$e^{-C_Tv}\geq n_0^{-\varepsilon/3}L^{-3C^2/\varepsilon}$. Hence
\[
 |T|\geq n_0^{1-\varepsilon/3}L^{-1-3C^2/\varepsilon}\geq\frac{n}{C_o}\,n^{-\varepsilon/3}
 (C_L\log x)^{-1-3C^2/\varepsilon}.
\]
Since $n\leq C_ox^{2+8\varepsilon}$ and $(2+8\varepsilon)/3<0.7$, the right side is at least
$x^{-\varepsilon}n$ for $x$ large.
\end{proof}

\begin{lemma}\label{hf-thick-orbits}
Let $\Phi$ be a Galois datum with group $A$, generator $a_0$ of $A_0$ and $n_0$ the order of $a_0$. Let
$L\geq1$, and let $T$ be a finite set of non-negative integers with no relations of length $L$ and
with $4Ls<n_0$ for all $s\in T$. Let $u_0,t,y\in A$ and put $T^\sharp=\{sa_0+u_0:s\in T\}$.
\begin{enumerate}
\item If the order of $y$ exceeds $4L$, then for every integer $i_0$ at most two integers
  $i\in[i_0,i_0+L)$ have $t+iy\in T^\sharp$.
\item If the order of $y$ is at most $4L$, the set $\{t+iy:i\in\mathbf Z\}$ meets $T^\sharp$ in at
  most one point.
\end{enumerate}
\end{lemma}
\leanref{Formalism/SharpRestriction/RamificationFrobenius/TotientGaloisClasses/FrameFacts.lean}{Formalism.HC.Frame.hostThickOrbit}

\begin{proof}
For (a), suppose that $a<b<c$ lie in $[i_0,i_0+L)$ and that $t+ay=s_1a_0+u_0$, $t+by=s_2a_0+u_0$ and
$t+cy=s_3a_0+u_0$ with $s_1,s_2,s_3\in T$. Then $(b-a)y=(s_2-s_1)a_0$ and $(c-a)y=(s_3-s_1)a_0$. If
$s_1=s_2$, then $(b-a)y=0$ with $0<b-a<L$, which is impossible since the order of $y$ exceeds $4L$;
in the same way $s_1\neq s_3$ and $s_2\neq s_3$. Moreover
\[
 \bigl((c-a)(s_2-s_1)-(b-a)(s_3-s_1)\bigr)a_0=(c-a)(b-a)y-(b-a)(c-a)y=0,
\]
so $n_0$ divides the integer $D=(c-a)(s_2-s_1)-(b-a)(s_3-s_1)$. Since $0<b-a,c-a<L$ and
$0\leq s_k<n_0/(4L)$, we have $|D|<n_0/2$, so $D=0$. This is a relation of length $L$ with
$i=b-a\neq j=c-a$, a contradiction.

For (b), suppose that $t+iy=s_1a_0+u_0$ and $t+i'y=s_2a_0+u_0$. Then $(i'-i)y=(s_2-s_1)a_0$. Let $m$
be the order of $h(y)$; it divides the order of $y$, so $m\leq4L$. Then
$h(m(s_2-s_1)a_0)=m(i'-i)h(y)=0$, and since $h$ is injective on $A_0$, $m(s_2-s_1)a_0=0$. So $n_0$
divides $m(s_2-s_1)$, whose absolute value is less than $4L\cdot n_0/(4L)=n_0$. Hence $s_1=s_2$ and
the two points coincide.
\end{proof}

\section{Hosts}\label{section-hf-hosts}

\begin{definition}\label{hf-adapted-ranks}
Let $s=(x,\varepsilon,J)$ be a scale (Definition~\ref{hf-windows}) and put $\delta_1=x^{-\varepsilon/8}$.
Let $R$ be a rank datum (Definition~\ref{th-ranks}) with integers $k,m,t$. Its data (a), with a
numbering of the group, form a Galois datum (Definition~\ref{hf-frame}), the \emph{Galois datum of}
$R$; let $n$ be the order of its group.
Put $f=2n$ and $r=k+m-1$. The rank datum is \emph{adapted to} $s$ if
\[
 x^{2+8\varepsilon}\leq n\leq2x^{2+8\varepsilon},\qquad
 \sqrt{x\delta_1}\leq m\leq2\sqrt{x\delta_1},\qquad
 x^{9/2+10\varepsilon}\leq k\leq2x^{9/2+10\varepsilon},
\]
and $r+t=8fr+4+T_xw$, where $T_x=\lceil x^{377}\rceil$ and
$w=8\bigl\lceil10r\bigl(f+f^2(k-1)^3x^{2\varepsilon}/m\bigr)/8\bigr\rceil$.
\end{definition}
\leanref{Development/AbelianVarietyTheory/Definitions/GlobalTwisting/PairRowIffQdRows.lean}{GlobalTheory.EX.RankVal}

\begin{proposition}\label{hf-adapted-ranks-exist}
Let $0<\varepsilon<1/52$. There is $x_0$ such that for every scale $s=(x,\varepsilon,J)$ with
$x\geq x_0$:
\begin{enumerate}
\item for every Galois datum $\Phi$ whose group has order $n$ with $x^{2+8\varepsilon}\leq n\leq
  2x^{2+8\varepsilon}$ there is a rank datum with Galois datum $\Phi$ adapted to $s$;
\item every rank datum adapted to $s$ has $t\leq x^{400}$.
\end{enumerate}
\end{proposition}
\leanref{Development/AbelianVarietyTheory/Proofs/GlobalInterpolation/ExistsRanksRankValFrame.lean}{GlobalTheory.EX.exists_ranks_rankVal_frame}
\leanref{Development/AutomorphicArithmetic/CharacterBoxes/EventuallyRankLeRpowOfLtOneDivFiftytwo.lean}{UpperTheory.Top.eventually_rank_le_rpow_of_lt_one_div_fiftytwo}

\begin{proof}
For (a), choose even $m$ and $k$ in their intervals, which have length at least $2$ for $x$ large,
and define $t$ by the equation of Definition~\ref{hf-adapted-ranks}. Since $r$ is odd and $8fr+4+T_xw$
is even, $t$ is odd. For $x$ large $t\geq8fr$ is much larger than $k$, $m$ and $r$, and
$m\leq2n\leq k$, so the inequalities of Definition~\ref{th-ranks} hold. For (b),
$t\leq8fr+4+T_xw$ with $f\leq4x^{2+8\varepsilon}$, $k\leq2x^{9/2+10\varepsilon}$,
$m\geq x^{1/2-\varepsilon/16}$ and $T_x\leq2x^{377}$; with $\varepsilon<1/52$ the exponent of $x$ in
the resulting bound is less than $400$.
\end{proof}

\begin{definition}\label{hf-host}
Let $R$ be a rank datum (Definition~\ref{th-ranks}) with Galois datum $\Phi$, group $A$ of order $n$, generator
$a_0$ of $A_0$, $n_0$ the order of $a_0$, and integers $k,m,t$; put $r=k+m-1$. A \emph{host} $H$
for $R$ consists of a host field $(q,d,Q_S,E,\gamma,j,\operatorname{art})$ on $\Phi$
(Definition~\ref{hf-framed-host-field}), with $N=Q_Sqd$ and $K=\mathbf Q(\sqrt{-d})$, and the
following data, subject to the stated conditions:
\begin{enumerate}
\item an integer $B\geq0$, the \emph{splitting bound}, such that the host field has splitting bound
  $B$ (Definition~\ref{hf-host-field-rows});
\item $q>4tr+1$ and $d>4tr+1$;
\item an embedding $\varphi_0\colon E\to\mathbf C$ with $\varphi_0(\sqrt{-d})=i\sqrt d$, and an
  ideal $\mathfrak b$ of $\mathcal O_E$ with the properties of
  Proposition~\ref{hf-host-field-properties}(f);
\item a finite set $\mathcal L$ of primes, the \emph{level primes}, prime to $Q_S$ and containing
  distinct primes $g_0,g_1$ and distinct primes $p,p'$ with $p\equiv p'\equiv1\pmod{15}$ and
  $p,p'\notin\{g_0,g_1\}$;
\item a nonzero ideal $\mathfrak c$ of $\mathcal O_K$, contained in no ideal $a\mathcal O_K$ with
  $a>1$ an integer, whose prime divisors have prime norms $\ell$ with $\left(\frac{-d}{\ell}\right)=1$,
  and whose norm is prime to $2$, $3$, $5$, $q$, $d$, $Q_S$ and the level primes;
\item a finite set $T_{\mathrm{top}}\subset A$ of \emph{top classes}, a finite set $T_2$ of
  non-negative integers, an element $u_0\in A$, an integer $L_w\geq2$ such that $T_2$ has no
  relations of length $L_w$ (Definition~\ref{hf-no-short-relations}) and $4L_ws<n_0$ for $s\in T_2$,
  and a real number $C_T$.
\end{enumerate}
The \emph{thick places} of $H$ are the elements $\tau_2(sa_0+u_0)$, $s\in T_2$, of $G_\Phi$.
\end{definition}
\leanref{Formalism/LocalRingCardinality/DataHost.lean}{Formalism.BasicNotions.Host}
\leanref{Formalism/LocalRingCardinality/DataHost.lean}{Formalism.BasicNotions.Host.thickPlaces}

\begin{lemma}\label{hf-host-extension}
Let $R$ be a rank datum with Galois datum $\Phi$, and let $(q,d,Q_S,E,\gamma,j,\operatorname{art})$ be a host field
on $\Phi$ with splitting bound $B$ (Definition~\ref{hf-host-field-rows}) and with $q,d>4tr+1$. Then
for every choice of the data (d), (e) and (f) of Definition~\ref{hf-host} that satisfies the
conditions there, there is a host for $R$ with this host field, splitting bound $B$ and these data.
\end{lemma}
\leanref{Development/AutomorphicArithmetic/CoprimeRanks/SecondSetNoRelation.lean}{HostParam.Seam.exists_host}
\leanref{Development/AutomorphicArithmetic/CoprimeRanks/AdmChi.lean}{UpperTheory.Top.admissible_ranks_of_CMDatum}

\begin{proof}
It remains to give (c). An embedding of $E$ into $\mathbf C$ sends $\sqrt{-d}$ to $\pm i\sqrt d$,
and composing with complex conjugation changes the sign, so $\varphi_0$ exists. The ideal
$\mathfrak b$ exists by Proposition~\ref{hf-host-field-properties}(f).
\end{proof}

\begin{definition}\label{hf-host-conditions}
Let $K$ be a level (Definition~\ref{wl-level}): a finite set of primitive Dirichlet characters that
contains the trivial character and is closed under complex conjugation, an integer $o_K\geq0$, and
either no exceptional character or an exceptional character $\psi_K\in K$, which is real and odd. Let $\Phi$ be a Galois datum with group $A$ of order $n$, index $o$ and
subgroup $A_0$, let $(q,d,Q_S,E,\gamma,j,\operatorname{art})$ be a host field on $\Phi$, let
$T\subset A$ be finite, and let $x,\varepsilon,C$ be real numbers. The host field satisfies the
\emph{host conditions} for $K$ at $(x,\varepsilon,C)$ with top classes $T$ if:
\begin{enumerate}
\item $o=o_K$;
\item every $\chi\in K$ is $\chi_\theta$ for some $\theta\in X^+_\Phi$;
\item if $K$ has an exceptional character $\psi_K$, then $\psi_K=\chi_{\theta_q}$; if not, then
  $\chi_{\theta_q}\notin K$;
\item for every $\gamma'\in A$,
  \[
   \Bigl|\#\{t\in T:t-\gamma'\in A_0\}-\lambda_{\gamma'}|T|\Bigr|\leq x^{-1+\varepsilon}|T|,
   \qquad\lambda_{\gamma'}=\begin{cases}4/(3o+1)&\gamma'\in A_0,\\3/(3o+1)&\gamma'\notin A_0;\end{cases}
  \]
\item for every $\theta\in X^-_\Phi$, and every $\theta\in X^+_\Phi$ with $\chi_\theta\notin K$,
  $|g_{\chi_\theta}(x)|\leq x^{1/2+\varepsilon}$ and
  $\left|\frac{L'}{L}(1,\chi_\theta)\right|\leq C\log x$.
\end{enumerate}
\end{definition}
\leanref{Formalism/SharpRestriction/RamificationFrobenius/TotientGaloisClasses/ExplicitFormulaErrorOfHostConditions.lean}{Formalism.HC.EAtH}
\leanref{Formalism/SharpRestriction/RamificationFrobenius/TotientGaloisClasses/ExplicitFormulaErrorOfHostConditions.lean}{Formalism.HC.Frame.Shares}

\begin{definition}\label{hf-good-orbit}
Let $x>1$ and $0<\varepsilon<1/52$ be real numbers, put $\delta_1=x^{-\varepsilon/8}$ and
$h_\varepsilon=2+8\varepsilon$, and put
\[
 K_\varepsilon=\bigl\lceil8(h_\varepsilon+2)/\varepsilon\bigr\rceil,\quad r_\varepsilon=K_\varepsilon+1,\quad
 u=\lceil1/\delta_1\rceil,\quad M_0=\bigl\lceil(h_\varepsilon+3)\,10^{6r_\varepsilon+1}\log x\bigr\rceil,
\]
\[
 G'=(M_0+5)(6r_\varepsilon+1)u,\qquad
 N_1=\bigl\lceil\max\bigl(6\exp(1)\,G'\delta_1,(h_\varepsilon+3)\log_2x\bigr)\bigr\rceil+1.
\]
Let $n'\geq1$ and $L_2\geq0$ be integers. For integers $i$ and $\ell\geq0$ let the \emph{window}
$W(i,\ell)\subset\{0,\ldots,n'-1\}$ be the set of residues modulo $n'$ of $i,i+1,\ldots,i+\ell-1$.
Put $S=\lfloor n'/((6r_\varepsilon+1)u)\rfloor$, and for integers $s\geq0$ and
$1\leq c\leq6r_\varepsilon+1$ let the \emph{cell} $I_{s,c}$ be the set of integers $i$ with
$s(6r_\varepsilon+1)u+(c-1)u\leq i<s(6r_\varepsilon+1)u+cu$. A set $P\subset\{0,\ldots,n'-1\}$ is a
\emph{tiled set of positions} for $(x,\varepsilon,L_2)$ if:
\begin{enumerate}
\item every window $W(i,L_2)$ with $0\leq i<n'$ contains at most $K_\varepsilon$ elements of $P$;
\item for every integer $0\leq s_0<S$ there is $0\leq j\leq M_0$ such that the \emph{tile} $s$, the
  residue of $s_0+j$ modulo $S$, is \emph{marked}: for odd $c$ the cell $I_{s,c}$ contains no element
  of $P$, and for even $c$ it contains exactly one element of $P$, and that element is even;
\item every window $W(i,G')$ with $0\leq i<n'$ contains at most $N_1$ elements of $P$.
\end{enumerate}
Now let $\Phi$ be a Galois datum with group $A$ (Definition~\ref{hf-frame}), let $y\in A$ have order
$n'$, and let $T\subset A$ and $t\in A$. The \emph{base point} of $t$ along $y$ is the element $b$ of
the coset $t+\langle y\rangle$ with the least number, and the \emph{positions} of $T$ along the orbit
of $t$ under $y$ are the integers $0\leq i<n'$ with $b+iy\in T$. The orbit of $t$ under $y$ is
\emph{tiled for} $T$ at $(x,\varepsilon,L_2)$ if these positions form a tiled set of positions for
$(x,\varepsilon,L_2)$.
\end{definition}
\leanref{Formalism/SubspaceCounting/GoodOrbitIntervalTileCondition.lean}{Formalism.G0.GoodOrbitIntervalTileCondition}
\leanref{Formalism/SharpRestriction/GrowthBounds/UnramifiedQuadraticSetting.lean}{Formalism.MainNotions.G.IsGoodOrbit}
\leanref{Formalism/LocalToGlobal/GlobalTheory/IsResidueClassZero.lean}{Formalism.L2.GlobalTheoryTheory.Arith.Loc.Good}
\leanref{Formalism/LocalToGlobal/GlobalTheory/IsResidueClassZero.lean}{Formalism.L2.GlobalTheoryTheory.Arith.bsH}

\begin{definition}\label{hf-good-host}
Let $s=(x,\varepsilon,J)$ be a scale (Definition~\ref{hf-windows}), and put
$\delta_1=x^{-\varepsilon/8}$, $b=\exp(x^\varepsilon)$ and $X'=\exp(x^{\varepsilon/3})$. Let $R$ be a rank datum with Galois datum $\Phi$, group $A$ of order $n$, index $o$,
$n_0$ the order of $a_0$, and integers $k,m,t$, and put $r=k+m-1$. Let $H$ be a host for $R$
(Definition~\ref{hf-host}) with the notation there, let $q_0$ be an integer, and let
$C_8,C_5,c_6,C_b$ be real numbers. For $Y>1$ put $z_Y=Y^{\varepsilon/5}$ and say that $Y$ is a
\emph{pair scale} for $H$ if $x^{1/J}\leq Y\leq x$, $e^{z_Y/2}<q\leq e^{z_Y}$,
$e^{z_Y\log^2Y/2}<d\leq e^{z_Y\log^2Y}$, $\left(\frac{-d}{q}\right)=1$,
$\left(\frac{-d}{p}\right)=1$ for every prime $p\leq z_Y$, and $\log N\leq2z_Y\log^2Y$. The host
$H$ is \emph{good at} $s$ with constants $(C_8,C_5,c_6,C_b)$ and auxiliary prime $q_0$ if:
\begin{enumerate}
\item there is a pair scale for $H$; $\delta_1n/2\leq|T_{\mathrm{top}}|\leq2\delta_1n$;
  $|T_2|\geq x^{-\varepsilon}n$; $4L_w|T_2|\leq n_0$;
  $|T_2|\geq\frac{n_0}{L_w}\exp(-C_T\sqrt{\log n_0\log L_w})$;
  $L_w=10\lceil405\log x\rceil+2\lceil30/\varepsilon\rceil+22$; $\mathfrak c=\mathcal O_K$;
  $B\geq\lfloor x^{\varepsilon/(5J)}\rfloor$; and $J\leq C_8$, $o\leq C_8$, $Q_S\leq C_8n$;
\item the level primes are $\mathcal L=\{g_0,g_1,p,p',q_0\}$; each $g_i$ satisfies
  $(kb)^8\leq g_i\leq2(kb)^8$, $g_i\equiv1\pmod N$, and $g_i=\pi\bar\pi$ for some
  $\pi\in\mathcal O_E\cap K$ with $\pi\equiv1\pmod{N\mathcal O_E}$; $p\equiv p'\equiv1\pmod{\operatorname{lcm}(15,N)}$ and
  $p,p'\leq N^{C_8}$; $q_0$ is a prime with $q_0\equiv1\pmod N$, $q_0\leq N^{C_8}$ and
  $q_0\notin\{p,p'\}$;
\item $p$ and $p'$ split completely in $E$; every level prime $\ell$ is prime to $N$ and has
  $\left(\frac{-d}{\ell}\right)=1$; and $\mathcal B_r(g_0)\cap\mathcal B_r(g_1)=\emptyset$, where
  $\mathcal B_r(g)$ is the set consisting of $g$ and the primes $\ell$ with
  $g^i<\ell^{4v_\ell(g^i-1)}$ for some $1\leq i\leq r$;
\item for $\chi\in\{\chi_{-q},\chi_{-d},\chi_{-q}\chi_{-d}\}$, $|g_\chi(x)|\leq x^{1/2+\varepsilon}$
  and $\left|\frac{L'}{L}(1,\chi)\right|\leq C_5\log x$; and
  $\bigl|\operatorname{Re}\frac{L'}{L}(1,\chi_{-d})\bigr|\leq x^{\varepsilon}$;
\item there is a pair scale $Y$ for $H$ such that, with $Q_1=e^{z_Y}$, $Q_2=e^{z_Y\log^2Y}$ and
  $Q_S'$ the odd part of $Q_S$: the characters $\chi_{-q}$ and $\chi_8\chi_{-q}$ have conductor at
  most $8Q_1$ and no zeros with $\operatorname{Re}s>1-1/400$, $|\operatorname{Im}s|\leq8Q_1$; the
  characters $\chi_{-d}$, $\chi_8\chi_{-d}$, $\chi_{-q}\chi_{-d}$ and $\chi_8\chi_{-q}\chi_{-d}$
  satisfy the same with $8Q_1Q_2$; and every real character $\chi_\theta$ of $E$ whose conductor is
  divisible by $d$, respectively by $q$ but not by $d$, satisfies the same with $8Q_S'Q_1Q_2$,
  respectively $8Q_S'Q_1$;
\item every character $\chi_\theta$ of $E$ satisfies $|g_{\chi_\theta}(y)|\leq
  y\exp(-c_6x^{\varepsilon/8})$ for all $y\geq X'/3$;
\item let $\mathcal P$ be the set of primes $p''<X'$ of inert class, that is, with $p''\nmid N$,
  $\left(\frac{-d}{p''}\right)=-1$ and $\epsilon_1(\operatorname{art}(p''))=-1$, for which the
  $A$-component $y_{p''}$ of $\operatorname{art}(p'')$ has order at least $x^{\varepsilon/4}$; let
  $L_2$ be twice the largest of the integers $\lceil4\log(kx^{26})/\log p\rceil$ over the primes $p$
  of inert class; then
  \[
   \sum_{p''\in\mathcal P}\frac{\log p''}{p''}\,\#\{t\in T_{\mathrm{top}}:\text{the orbit of }t
   \text{ under }y_{p''}\text{ is not tiled for }T_{\mathrm{top}}\text{ at }(x,\varepsilon,L_2)\}
   \leq4C_bx^{-2}\delta_1n\sum_{p''\in\mathcal P}\frac{\log p''}{p''}
  \]
  (Definition~\ref{hf-good-orbit}).
\end{enumerate}
\end{definition}
\leanref{Development/AbelianVarietyTheory/Definitions/GlobalTwisting/ThinScalesIffHostBound.lean}{GlobalTheory.EX.HostGood}
\leanref{Development/AbelianVarietyTheory/Definitions/GlobalTwisting/PairRowIffQdRows.lean}{GlobalTheory.EX.HostVal}
\leanref{Development/AbelianVarietyTheory/Definitions/GlobalTwisting/PairRowIffQdRows.lean}{GlobalTheory.EX.HostArith}
\leanref{Development/AbelianVarietyTheory/Definitions/GlobalTwisting/PairRowIffQdRows.lean}{GlobalTheory.EX.PairRow}
\leanref{Formalism/SharpRestriction/InertJacobiLevels/AsmIface.lean}{Formalism.MainNotions.Inhab.QDq}
\leanref{Formalism/SharpRestriction/InertJacobiLevels/AsmIface.lean}{Formalism.MainNotions.Inhab.Boxq}
\leanref{Development/AbelianVarietyTheory/Definitions/GlobalTwisting/ThinScalesIffHostBound.lean}{GlobalTheory.EX.BoxW}
\leanref{Development/AbelianVarietyTheory/Definitions/GlobalTwisting/ThinScalesIffHostBound.lean}{GlobalTheory.EX.ThinOK}
\leanref{Development/AbelianVarietyTheory/Definitions/GlobalTwisting/ThinScalesIffHostBound.lean}{GlobalTheory.EX.BadOK}
\leanref{Formalism/CyclotomicNorms/StatementI.lean}{Formalism.CyclotomicNorms.BB}

\begin{theorem}\label{hf-norm-primes-short}
There are $A,c,c'>0$ with the following property. Let $d\geq100$ be a prime with $d\equiv3\pmod4$,
let $M\geq3$ be an integer, and let $X\geq1$ and $h$ be real numbers with
$\log(4dM^2)\leq c\sqrt{\log X}$ and $\exp(-c\sqrt{\log X})\leq h\leq1$. Then the number of primes
$g$ with $X<g\leq X(1+h)$ for which $4g=a^2+dc^2$ with integers $a,c$ such that $M\mid c$ and
$2M\mid a-2-c$ is at least
\[
 \frac{c'hX}{(4dM^2)^A\log X}.
\]
\end{theorem}
\leanref{Development/AbelianVarietyTheory/Proofs/GlobalInterpolation/NormPrimesImagQuadShortIntervalLowerBoundStatement.lean}{GlobalTheory.T.normPrimesImagQuad_short_interval_lower_bound}

With $K=\mathbf Q(\sqrt{-d})$ and $\pi=(a+c\sqrt{-d})/2$, the condition on $g$ says that
$g=\pi\bar\pi$ with $\pi\in\mathcal O_K$ and $\pi\equiv1\pmod{M\mathcal O_K}$, that is, $g$ splits
completely in the ray class field of $K$ modulo $M$. The bound follows from the prime ideal theorem
for the Hecke $L$-functions of $K$ with the classical zero-free region, in which a possible
exceptional zero is controlled by Stark's bound \cite{LO77,Sta74}.

\begin{lemma}\label{hf-exceptional-primes}
For an integer $r\geq1$ and an integer $g\geq2$ let $\mathcal B_r(g)$ be the set consisting of $g$
and the primes $\ell$ with $g^i<\ell^{4v_\ell(g^i-1)}$ for some $1\leq i\leq r$. Let $Y\geq2$, and
let $g_1$ be an integer with $Y\leq g_1\leq2Y$. Then at most
$(3r+1)\bigl(1+2r^2(Y^{3/4}+1)\bigr)$ integers $g$ with $Y\leq g\leq2Y$ satisfy
$\mathcal B_r(g)\cap\mathcal B_r(g_1)\neq\emptyset$.
\end{lemma}
\leanref{Development/AbelianVarietyTheory/Proofs/FoundationProofs/DivLeRpowThreeQuarters.lean}{GlobalTheory.choice_count}

\begin{proof}
First, $|\mathcal B_r(g_1)|\leq3r+1$. Indeed, for each $i$ at most three primes $\ell$ satisfy
$g_1^i<\ell^{4v_\ell(g_1^i-1)}$: for four of them the product of the $\ell^{v_\ell(g_1^i-1)}$ would
divide $g_1^i-1$ and exceed $g_1^i$.

Next, fix $\ell$; we show that at most $1+2r^2(Y^{3/4}+1)$ integers $g\in[Y,2Y]$ have
$\ell\in\mathcal B_r(g)$. One of them may be $g=\ell$. Otherwise $\ell$ is a prime with
$g^i<\ell^{4v}$ for some $i\leq r$, where $v=v_\ell(g^i-1)$. Let $w_i$ be the least integer with
$\ell^{4w_i}>Y^i$. Then $v\geq w_i$, so $g^i\equiv1\pmod{\ell^{w_i}}$, and $\ell^{w_i}>Y^{1/4}$.
There are at most $2i$ residues $x$ modulo $\ell^{w_i}$ with $x^i\equiv1$, since
$(\mathbf Z/\ell^{w}\mathbf Z)^\times$ is cyclic for odd $\ell$ and is the product of a cyclic group
with a group of order $2$ for $\ell=2$. So at most $2i(Y/\ell^{w_i}+1)\leq2r(Y^{3/4}+1)$ integers
$g\in[Y,2Y]$ arise from $i$, and summing over $i\leq r$ gives the bound.

Every $g$ counted in the statement has an element $\ell\in\mathcal B_r(g)\cap\mathcal B_r(g_1)$.
Multiplying the two bounds gives the claim.
\end{proof}

\begin{lemma}\label{hf-norm-primes}
There is $C_1\geq1$ such that for every $0<\varepsilon<1/52$ there is $x_0$ with the following
property. Let $s=(x,\varepsilon,J)$ be a scale with $x\geq x_0$ and $b=\exp(x^\varepsilon)$, let $R$ be
a rank datum adapted to $s$ (Definition~\ref{hf-adapted-ranks}) with integers $k,m$ and $r=k+m-1$, let
$d\geq100$ be a prime with $d\equiv3\pmod4$, put $K=\mathbf Q(\sqrt{-d})$, and let $N\geq1$ be an
integer with $d\mid N$ and $\log N\leq3x^{\varepsilon/5}\log^2x$. Then there are primes $g_0\neq g_1$, $p\neq p'$ and $q_0$ such
that:
\begin{enumerate}
\item $(kb)^8\leq g_i\leq2(kb)^8$, and $4g_i=a^2+dc^2$ with integers $a,c$ such that $N\mid c$ and
  $2N\mid a-2-c$; equivalently $g_i=\pi\bar\pi$ with $\pi=(a+c\sqrt{-d})/2\in\mathcal O_K$ and
  $\pi\equiv1\pmod{N\mathcal O_K}$;
\item $\mathcal B_r(g_0)\cap\mathcal B_r(g_1)=\emptyset$, where $\mathcal B_r(g)$ is the set
  consisting of $g$ and the primes $\ell$ with $g^i<\ell^{4v_\ell(g^i-1)}$ for some $1\leq i\leq r$;
\item $p\equiv p'\equiv1\pmod{15N}$, $q_0\equiv1\pmod N$, and $p,p',q_0\leq N^{C_1}$;
\item $q_0\notin\{p,p'\}$ and $p,p'\notin\{g_0,g_1\}$.
\end{enumerate}
\end{lemma}
\leanref{Development/AutomorphicArithmetic/ShortNorms/levelData.lean}{UpperTheory.T3.levelData}
\leanref{Development/AutomorphicArithmetic/ShortNorms/Primes.lean}{UpperTheory.T3.shortNorm}
\leanref{Development/AutomorphicArithmetic/ShortNorms/Primes.lean}{UpperTheory.T3.smallOne}

\begin{proof}
Let $A'\geq1$ be such that for every integer $M\geq2$ the least prime $p\equiv1\pmod M$ is at most
$M^{A'}$; it exists by Linnik's theorem (Theorem~\ref{lit-linnik} in
Chapter~\ref{chapter-literature}). Put $C_1=(1+2A'+(2A'+2)A')A'$, and let $x$ be large.

\emph{Step 1: many norm primes.} Put $X=(kb)^8$. Since $k\geq1$, $\log X\geq8x^\varepsilon$. Apply
Theorem~\ref{hf-norm-primes-short} with $M=N$ and $h=1$; here $N\geq d\geq100$. Its hypothesis holds
for $x$ large, because $\log(4dN^2)\leq\log4+3\log N\leq10x^{2\varepsilon/5}$ while
$c\sqrt{\log X}\geq c\sqrt8\,x^{\varepsilon/2}$. It gives at least
$c'X/((4dN^2)^A\log X)\geq c'X\exp(-10Ax^{2\varepsilon/5})/\log X\geq X^{9/10}$ primes $g$ with
$X<g\leq2X$ satisfying the condition of (a).

\emph{Step 2: a disjoint pair.} Since $R$ is adapted to $s$, $r\leq4x^{9/2+10\varepsilon}\leq x^5$
for $x$ large, so $20r^3\leq20x^{15}<\exp(x^\varepsilon)\leq X^{3/20}$. Then
$(3r+1)(1+2r^2(X^{3/4}+1))\leq20r^3X^{3/4}<X^{9/10}$. Let $g_0$ be one of the primes of Step 1. By
Lemma~\ref{hf-exceptional-primes}, fewer than $X^{9/10}$ of these primes $g$ have
$\mathcal B_r(g)\cap\mathcal B_r(g_0)\neq\emptyset$, so there is a prime $g_1$ among them with
$\mathcal B_r(g_1)\cap\mathcal B_r(g_0)=\emptyset$. Since $g_0\in\mathcal B_r(g_0)$, $g_1\neq g_0$.
This gives (a) and (b).

\emph{Step 3: the small primes.} Let $p$ be the least prime with $p\equiv1\pmod{15N}$, $p'$ the least
prime with $p'\equiv1\pmod{15Np}$, and $q_0$ the least prime with $q_0\equiv1\pmod{Npp'}$. Then
$p'\equiv1\pmod p$, so $p'\neq p$, and $q_0\equiv1$ modulo $p$ and $p'$, so $q_0\notin\{p,p'\}$.
Also $p'\equiv1\pmod{15N}$ and $q_0\equiv1\pmod N$. Since $N\geq15$, $p\leq(15N)^{A'}\leq N^{2A'}$,
$p'\leq(15Np)^{A'}\leq N^{(2A'+2)A'}$ and $q_0\leq(Npp')^{A'}\leq N^{C_1}$. This gives (c) and the
first part of (d).

\emph{Step 4: separation.} Since $\log N\leq3x^{\varepsilon/5}\log^2x$,
$C_1\log N<x^\varepsilon$ for $x$ large, so $p,p'\leq N^{C_1}<b\leq X\leq g_i$. Hence
$p,p'\notin\{g_0,g_1\}$.
\end{proof}

\begin{theorem}\label{hf-pnt-character}
For a Dirichlet character $\chi$ and $u>0$ put $\psi(u,\chi)=\sum_{n\leq u}\Lambda(n)\chi(n)$, and
let $\delta_\chi=1$ if $\chi$ is trivial and $\delta_\chi=0$ otherwise. For every $c_1>0$ there are
$C,c>0$ with the following property. Let $\chi$ be a primitive Dirichlet character of conductor $N$,
let $\lambda\geq1$ with $\log N\leq\lambda$, and suppose that $L(\beta,\chi)\neq0$ for
$1-c_1<\beta<1$ if $\chi$ is real and nontrivial. Then for every $u\geq e$,
\[
 |\psi(u,\chi)-\delta_\chi u|\leq Cu(\log u+\lambda)^2\exp\Bigl(-\frac{c\log u}{\lambda+\sqrt{\log u}}\Bigr).
\]
\end{theorem}
\leanref{Development/AutomorphicArithmetic/MeanSquareBounds/PsiThin.lean}{QH.psi_thin}

This is the prime number theorem for $\chi$ obtained from the explicit formula for $\psi(u,\chi)$
and the classical zero-free region, in which only a real zero of a real character can occur
\cite[Chapters 14, 19 and 20]{Dav00}.

\begin{lemma}\label{hf-chebyshev-integral}
Let $\chi$ be a primitive Dirichlet character and $y\geq1$. Then, with the notation of
Theorem~\ref{hf-pnt-character},
\[
 g_\chi(y)=-y\int_y^\infty\frac{\psi(u,\chi)-\delta_\chi u}{u^2}\,du.
\]
\end{lemma}
\leanref{Development/AutomorphicArithmetic/MeanSquareBounds/GIntegral.lean}{QH.g_integral}

\begin{proof}
For $X>y$, $\int_y^X\psi(u,\chi)u^{-2}\,du=\psi(y,\chi)/y+\sum_{y<n\leq X}\Lambda(n)\chi(n)/n-
\psi(X,\chi)/X$. If $\chi\neq1$, let $X\to\infty$: $\psi(X,\chi)/X\to0$ and
$\sum_{n>y}\Lambda(n)\chi(n)/n=-\frac{L'}{L}(1,\chi)-\sum_{n\leq y}\Lambda(n)\chi(n)/n$; multiplying
by $-y$ gives $\sum_{n\leq y}\Lambda(n)\chi(n)(y/n-1)+y\frac{L'}{L}(1,\chi)=g_\chi(y)$. If $\chi=1$,
subtract $\int_y^Xu^{-1}\,du=\log(X/y)$ and use $\psi(X)/X\to1$ and
$\sum_{n\leq X}\Lambda(n)/n=\log X-\gamma+o(1)$; the limit is
$\psi(y)/y-\sum_{n\leq y}\Lambda(n)/n+\log y-1-\gamma$, and multiplying by $-y$ gives $G(y)$.
\end{proof}

\begin{lemma}\label{hf-host-chebyshev}
There is $c>0$ such that for every $B_0$ and $0<\varepsilon<1/52$ there is $x_0$ with the
following property. Let $x\geq x_0$, let $\Phi$ be a Galois datum and $(q,d,Q_S,E,\gamma,j,\operatorname{art})$
a host field on $\Phi$ with $N=Q_Sqd$, such that the odd part of $Q_S$ is at most $B_0$,
$\log N\leq2x^{\varepsilon/5}\log^2x$, and $L(\beta,\chi_\theta)\neq0$ for every real
$\beta>1-1/400$ and every real character $\chi_\theta$ of $E$ whose conductor is divisible by $q$
or $d$. Then every character $\chi_\theta$ of $E$ satisfies
$|g_{\chi_\theta}(y)|\leq y\exp(-cx^{\varepsilon/8})$ for all $y\geq\exp(x^{\varepsilon/3})/3$.
\end{lemma}
\leanref{Development/AutomorphicArithmetic/MeanSquareBounds/GaussSumFunctionEq.lean}{QH.thinH}
\leanref{Development/AutomorphicArithmetic/CharacterBoxes/PairGood.lean}{UpperTheory.Top.ThinH}

\begin{proof}
We prove the statement with $c=1$. Let $Q_S'$ be the odd part of $Q_S$.

\emph{Step 1: small moduli.} There are finitely many nontrivial Dirichlet characters of modulus at
most $8B_0$. For each of them $L(\beta,\chi)$ is continuous in $\beta$ near $1$ and
$L(1,\chi)\neq0$. So there is $c_0>0$ such that none of them has a zero $\beta$ with
$1-c_0<\beta<1$. Put $c_1=\min(c_0,1/400)$.

\emph{Step 2: real characters of $E$.} Let $\chi$ be a real nontrivial character of $E$, primitive
of conductor $f\mid N$. If $q\mid f$ or $d\mid f$, then $L(\beta,\chi)\neq0$ for $1-c_1<\beta<1$ by
hypothesis. Otherwise, since $\chi$ is real and primitive, $f$ is $2^a$ times an odd squarefree
integer with $a\leq3$. Every prime factor of $f$ divides $N=Q_Sqd$ and is neither $q$ nor $d$, so
$f\mid8Q_S'$ and $f\leq8B_0$, and Step 1 applies. So the hypothesis of
Theorem~\ref{hf-pnt-character} holds for every character of $E$ with this $c_1$; let $C,c$ be the
constants it gives.

\emph{Step 3: the integral.} Let $\chi$ be a character of $E$ of conductor $f$, put
$\lambda=\max(1,\log f)$, and let $y\geq\exp(x^{\varepsilon/3})/3$. Then
$\lambda\leq3x^{\varepsilon/5}\log^2x$ for $x$ large, because $\log f\leq\log N$. By
Lemma~\ref{hf-chebyshev-integral} and Theorem~\ref{hf-pnt-character}, with $v=\log u$,
\[
 |g_\chi(y)|\leq Cy\int_{\log y}^\infty(v+\lambda)^2\exp\Bigl(-\frac{cv}{\lambda+\sqrt v}\Bigr)\,dv.
\]

\emph{Step 4: the exponent.} On the range of integration $v\geq x^{\varepsilon/3}/2$ for $x$ large.
If $v\geq\lambda^2$, then $cv/(\lambda+\sqrt v)\geq c\sqrt v/2\geq cx^{\varepsilon/6}/3$; if
$v<\lambda^2$, then $cv/(\lambda+\sqrt v)\geq cv/(2\lambda)\geq cx^{2\varepsilon/15}/(12\log^2x)$.
Since $1/6>2/15>1/8$, half of the exponent is at least $2x^{\varepsilon/8}$ for $x$ large. The
integral of $(v+\lambda)^2$ against the exponential of the other half is $O(\lambda^5+1)$, which is
at most $\exp(x^{\varepsilon/8})/C$ for $x$ large. Hence $|g_\chi(y)|\leq y\exp(-x^{\varepsilon/8})$.
\end{proof}

\begin{lemma}\label{hf-orbit-failure}
For all $0<\varepsilon<1/52$, $C_o\geq1$ and $L_b\geq0$ there are $C_b>0$ and $x_0$ with the
following property. Let $x\geq x_0$ and $\delta_1=x^{-\varepsilon/8}$, let $n'$ be an integer with
$x^{\varepsilon/4}\leq n'\leq C_ox^{2+8\varepsilon}$, let $0\leq L_2\leq L_b$ and $0\leq i_0<n'$ be
integers, and let $P$ be a random subset of $\{0,\ldots,n'-1\}\setminus\{i_0\}$ that contains each
$i$ independently with a probability $p_i$ such that $3\delta_1/4\leq p_i\leq\delta_1$. Then the
probability that $P\cup\{i_0\}$ is not a tiled set of positions for $(x,\varepsilon,L_2)$
(Definition~\ref{hf-good-orbit}) is at most $C_bx^{-2}$.
\end{lemma}
\leanref{Development/AutomorphicArithmetic/ThickFrames/NotGoodLe.lean}{UpperTheory.BernoulliSumsFrobeniusOrbits.not_good_le}
\leanref{Development/AutomorphicArithmetic/ThickFrames/orbitTerms.lean}{UpperTheory.BernoulliSumsFrobeniusOrbits.orbit_terms}

\begin{proof}
Use the notation $K_\varepsilon,r_\varepsilon,u,M_0,G',N_1,S$ of Definition~\ref{hf-good-orbit}, and
write $r=r_\varepsilon$. Let $x$ be large.

\emph{Step 1.} Since $M_0\ll\log x$, $u\leq x^{\varepsilon/8}+1$ and $n'\geq x^{\varepsilon/4}$, we
have $(M_0+1)(6r+1)u\leq n'$, so $S\geq M_0+1$.

\emph{Step 2: condition (a).} If (a) fails for $P\cup\{i_0\}$, then some window $W(i,L_2)$ contains
at least $K_\varepsilon$ elements of $P$. For a fixed window this has probability at most
$\binom{L_2}{K_\varepsilon}\delta_1^{K_\varepsilon}$, since each of its at most $L_2$ elements lies in
$P$ with probability at most $\delta_1$. So (a) fails with probability at most
$n'\binom{L_2}{K_\varepsilon}\delta_1^{K_\varepsilon}$.

\emph{Step 3: condition (c).} In the same way, (c) fails with probability at most
$n'\Pr(\beta\geq N_1-1)$, where $\beta$ is a binomial random variable with $G'$ trials and success
probability $\delta_1$.

\emph{Step 4: condition (b).} The tiles $[s(6r+1)u,(s+1)(6r+1)u)$, $0\leq s<S$, are disjoint subsets
of $\{0,\ldots,n'-1\}$, and whether tile $s$ is marked depends only on its intersection with $P$. A
tile not containing $i_0$ is marked with probability at least
\[
 \mu=\bigl((1-\delta_1)^u\bigr)^{3r+1}\Bigl(\lfloor u/2\rfloor\tfrac34\delta_1(1-\delta_1)^{u-1}\Bigr)^{3r}:
\]
its $3r+1$ cells of odd index must miss $P$, and each of its $3r$ cells of even index must meet $P$
in one of its at least $\lfloor u/2\rfloor$ even integers only. For fixed $s_0$, the $M_0+1$ tiles
with indices $s_0,\ldots,s_0+M_0$ modulo $S$ are distinct by Step 1, and at most one of them contains
$i_0$. So (b) fails with probability at most $S(1-\mu)^{M_0}\leq n'(1-\mu)^{M_0}$.

\emph{Step 5: the bounds.} Since $u=\lceil1/\delta_1\rceil$ and $\delta_1\leq1/100$,
$\mu\geq10^{-(6r+1)}$, so $\mu M_0\geq(h_\varepsilon+3)\log x$ and
$n'(1-\mu)^{M_0}\leq C_ox^{2+8\varepsilon}x^{-5-8\varepsilon}\leq C_ox^{-3}$. Since
$\varepsilon K_\varepsilon/8\geq h_\varepsilon+2=4+8\varepsilon$,
$n'\binom{L_2}{K_\varepsilon}\delta_1^{K_\varepsilon}\leq C_oL_b^{K_\varepsilon}x^{-2}$. With
$m=N_1-1\geq6\exp(1)G'\delta_1$,
$\Pr(\beta\geq m)\leq\sum_{j\geq m}\binom{G'}{j}\delta_1^j\leq\sum_{j\geq m}(\exp(1)G'\delta_1/j)^j
\leq\sum_{j\geq m}6^{-j}\leq2^{-m}$, and $2^{-m}\leq x^{-(h_\varepsilon+3)}$ because
$m\geq(h_\varepsilon+3)\log_2x$; so $n'\Pr(\beta\geq N_1-1)\leq C_ox^{-3}$. The lemma follows
with $C_b=C_oL_b^{K_\varepsilon}+2C_o$.
\end{proof}

\begin{lemma}\label{hf-random-selection}
Let $\mathcal A$ be a finite set divided into classes, let $0\leq\delta\leq1$, let
$p\colon\mathcal A\to[0,\delta]$, and let $T$ be the random subset of $\mathcal A$ containing each
$a$ independently with probability $p(a)$. Let $D>0$ with $2\delta|\mathcal A|\leq D^2$ and let
$\beta\geq0$. For $i$ in a finite set $I$ let $\omega_i\geq0$ and let $\mathrm{BAD}_i(t,T')$ be a
condition on $t\in\mathcal A$ and $T'\subset\mathcal A$, such that the probability that $t\in T$
and $\mathrm{BAD}_i(t,T)$ holds is at most $\beta\delta$ for every $i\in I$ and $t\in\mathcal A$.
Then there is a subset $T'\subset\mathcal A$ such that
$\bigl|\#(T'\cap c)-\sum_{a\in c}p(a)\bigr|\leq D$ for every class $c$, and
\[
 \sum_{i\in I}\omega_i\,\#\{t\in T':\mathrm{BAD}_i(t,T')\}\leq4\beta\delta|\mathcal A|\sum_{i\in I}\omega_i.
\]
\end{lemma}
\leanref{Development/AutomorphicArithmetic/ThickFrames/select.lean}{UpperTheory.BernoulliSumsFrobeniusOrbits.select}

\begin{proof}
For a class $c$ let $Y_c=\#(T\cap c)-\sum_{a\in c}p(a)$. It has mean $0$ and variance
$\sum_{a\in c}p(a)(1-p(a))\leq\delta|c|$, so the expectation of $\sum_cY_c^2$ is at most
$\delta|\mathcal A|\leq D^2/2$. Let $B=\sum_i\omega_i\#\{t\in T:\mathrm{BAD}_i(t,T)\}$. Its expectation
is $\sum_i\omega_i\sum_t\Pr(t\in T,\ \mathrm{BAD}_i(t,T))\leq m$, where
$m=\beta\delta|\mathcal A|\sum_i\omega_i$. If $m=0$, then $B=0$ for every value of $T$ of positive
probability; otherwise the expectation of $\sum_cY_c^2/D^2+B/(4m)$ is at most $3/4$. In both cases
some value $T'$ of $T$ of positive probability has $\sum_cY_c^2\leq D^2$ and $B\leq4m$.
\end{proof}

\begin{lemma}\label{hf-top-classes}
For every $0<\varepsilon<1/52$ and $C_8$ there are $C_b>0$ and $x_0$ with the following property.
Let $s=(x,\varepsilon,J)$ be a scale with $x\geq x_0$, let $R$ be a rank datum adapted to $s$ whose Galois datum
has group $A$ of order $n$ and index $o$, and let $H$ be a host for $R$ with $J\leq C_8$, $o\leq C_8$,
$Q_S\leq C_8n$ and a pair scale (Definition~\ref{hf-good-host}). Then there is $T\subset A$ with
$\delta_1n/2\leq|T|\leq2\delta_1n$ that satisfies condition (d) of
Definition~\ref{hf-host-conditions} at $(x,\varepsilon)$, and such that $H$ with $T_{\mathrm{top}}$
replaced by $T$ satisfies condition (g) of Definition~\ref{hf-good-host}.
\end{lemma}
\leanref{Development/AutomorphicArithmetic/ThickFrames/topsChoiceH.lean}{UpperTheory.BernoulliSumsFrobeniusOrbits.topsChoiceH}

\begin{proof}
Put $C=\max(C_8,2)$ and $L_b=2(620C/\varepsilon+1)$, and let $C_b$ be the constant of
Lemma~\ref{hf-orbit-failure} for $\varepsilon$, $C_o=C$ and $L_b$. Let $x$ be large, let $n_0=n/o$
be the order of $A_0$, and let $L_2$ be as in condition (g) of Definition~\ref{hf-good-host}.

\emph{Step 1: primes of inert class are large.} Let $Y$ be a pair scale for $H$. Then
$\left(\frac{-d}{p}\right)=1$ for every prime $p\leq Y^{\varepsilon/5}$, and
$Y^{\varepsilon/5}\geq x^{\varepsilon/(5J)}\geq x^{\varepsilon/(5C)}$. A prime of inert class has
$\left(\frac{-d}{p}\right)=-1$, so it exceeds $x^{\varepsilon/(5C)}$.

\emph{Step 2: $L_2\leq L_b$.} Since $k\leq2x^{9/2+10\varepsilon}$ and $\varepsilon<1/52$,
$4\log(kx^{26})\leq124\log x$ for $x\geq16$. For $p$ of inert class $\log p>\varepsilon\log x/(5C)$
by Step 1, so $\lceil4\log(kx^{26})/\log p\rceil\leq620C/\varepsilon+1$.

\emph{Step 3: the random set.} Let $T$ be the random subset of $A$ that contains each $a$
independently with probability $p(a)=\delta_1$ if $a\in A_0$ and $p(a)=3\delta_1/4$ otherwise.

\emph{Step 4: one orbit.} Let $p''\in\mathcal P$, let $y=y_{p''}$, of order $n'$, and let $t\in A$.
Then $x^{\varepsilon/4}\leq n'\leq n\leq2x^{2+8\varepsilon}$. Let $b$ be the base point of $t$ along
$y$ and $i_0$ the integer with $b+i_0y=t$. The map $i\mapsto b+iy$ is injective on
$\{0,\ldots,n'-1\}$, so the positions of $T$ along the orbit of $t$ form a random set containing
each $i$ independently with probability $p(b+iy)\in[3\delta_1/4,\delta_1]$, and $t\in T$ exactly
when $i_0$ is a position. By Lemma~\ref{hf-orbit-failure}, the probability that $t\in T$ and the
orbit of $t$ under $y$ is not tiled for $T$ at $(x,\varepsilon,L_2)$ is at most
$p(t)C_bx^{-2}\leq C_bx^{-2}\delta_1$.

\emph{Step 5: the choice.} Apply Lemma~\ref{hf-random-selection} to $\mathcal A=A$ with the cosets
of $A_0$ as classes, $\delta=\delta_1$, $D=\sqrt{2\delta_1n}$, $\beta=C_bx^{-2}$, $I=\mathcal P$,
$\omega_{p''}=\log p''/p''$, and $\mathrm{BAD}_{p''}(t,T')$ the condition that the orbit of $t$ under
$y_{p''}$ is not tiled for $T'$ at $(x,\varepsilon,L_2)$. By Step 4 it gives a set $T'$ satisfying
condition (g), with $|\#(T'\cap(\gamma+A_0))-p_\gamma\delta_1n_0|\leq D$ for every $\gamma\in A$,
where $p_\gamma=1$ if $\gamma\in A_0$ and $p_\gamma=3/4$ otherwise.

\emph{Step 6: the shares.} Since $o\leq C$ and $\delta_1n\geq x^{2+7\varepsilon}/2$, for $x$ large
$32o^4x^2\leq\delta_1n$, which is equivalent to $4oD\leq x^{-1}\delta_1n_0$. Summing over the $o$
cosets, $|T'|=\frac{3o+1}4\delta_1n_0+\theta oD$ with $|\theta|\leq1$. Since
$\frac{3o+1}{4}n_0$ lies between $\frac34n$ and $n$, this gives $\delta_1n/2\leq|T'|\leq2\delta_1n$
and $|T'|\geq\frac34\delta_1n_0$. With $\lambda_\gamma=4p_\gamma/(3o+1)$ as in condition (d) of
Definition~\ref{hf-host-conditions}, $\lambda_\gamma|T'|=p_\gamma\delta_1n_0+\theta'D$ with
$|\theta'|\leq2$, so
\[
 \bigl|\#(T'\cap(\gamma+A_0))-\lambda_\gamma|T'|\bigr|\leq3D\leq\tfrac34x^{-1}\delta_1n_0\leq
 x^{-1}|T'|\leq x^{-1+\varepsilon}|T'|.
\]
This is condition (d). Take $T=T'$.
\end{proof}

\begin{theorem}\label{hf-good-host-exists}
Let $0<\varepsilon<1/52$ and $C_0\geq1$. There are real numbers $C_T$, $C_8\geq C_0$, $c_6>0$,
$C_b>0$ and $x_0$ with the following property. Let $s=(x,\varepsilon,J)$ be a scale with $x\geq x_0$
and $J\leq C_0$, let $R$ be a rank datum adapted to $s$ (Definition~\ref{hf-adapted-ranks}) whose Galois datum
$\Phi$ has $2$-power exponent $e$, group of order $n$ and index $o\leq C_0$, and let
$(q,d,Q_S,E,\gamma,j,\operatorname{art})$ be a host field on $\Phi$ with $Q_S\leq C_0n$. Let
$1\leq M\leq C_0$ be an integer with $Q_S\mid\operatorname{lcm}(2^{e+2},M)$, and let $Y$ with
$x^{1/J}\leq Y\leq x$ be such that $(q,d)$ is good for $(M,\varepsilon,Y,x)$
(Definition~\ref{hf-good-pair}) and the host field has splitting bound
$\lfloor Y^{\varepsilon/5}\rfloor$ (Definition~\ref{hf-host-field-rows}). Then there are a host $H$
for $R$ with this host field and constant $C_T$, and a prime $q_0$, such that $H$ is good at $s$
with constants $(C_8,4,c_6,C_b)$ and auxiliary prime $q_0$ (Definition~\ref{hf-good-host}), and its
top classes satisfy condition (d) of Definition~\ref{hf-host-conditions} at $(x,\varepsilon)$.
\end{theorem}
\leanref{Development/AutomorphicArithmetic/CharacterBoxes/HostOfField.lean}{UpperTheory.Top.hostOfField}
\leanref{Development/AutomorphicArithmetic/ShortNorms/hostOfFieldOf.lean}{UpperTheory.T3.hostOfFieldOf}
\leanref{Development/AutomorphicArithmetic/CoprimeRanks/AdmChi.lean}{UpperTheory.Top.admissible_ranks_of_CMDatum}

\begin{proof}
Let $C_1$ be the constant of Lemma~\ref{hf-norm-primes}, put $C_8=\max(C_0,C_1)$, let $c_6$ be the
constant of Lemma~\ref{hf-host-chebyshev}, let $C_b$ be the constant of Lemma~\ref{hf-top-classes}
for $\varepsilon$ and $C_0$, and let $C_T$ be the constant of Corollary~\ref{hf-behrend-set-at-scale}.
Take $x$ large. Write $N=Q_Sqd$, $z=Y^{\varepsilon/5}$ and $L_w=10\lceil405\log x\rceil+
2\lceil30/\varepsilon\rceil+22$.

Since $Y\geq x^{1/C_0}$, $\log q>z/2\geq x^{\varepsilon/(5C_0)}/2$, while $t\leq x^{400}$ by
Proposition~\ref{hf-adapted-ranks-exist} and $r\leq4x^{9/2+10\varepsilon}$; so $q,d>4tr+1$ and
$d\geq100$. Since $Q_S\leq2C_0x^{2+8\varepsilon}$, $\log q\leq z$ and $\log d\leq z\log^2Y$, we get
$\log N\leq2z\log^2Y$; with Definition~\ref{hf-admissible-pair} this shows that $Y$ is a pair scale.

The good pair gives the remaining analytic conditions on $q$ and $d$. Taking $\xi=1$ in
Definition~\ref{hf-good-pair}(b) at $y=x$ gives condition (d) of Definition~\ref{hf-good-host} with
$C_5=4$, and $|\operatorname{Re}\frac{L'}{L}(1,\chi_{-d})|\leq4\log x\leq x^{\varepsilon}$. Taking
$\rho\in\{1,\chi_8\}$ in (c) gives the zero-free boxes for the six quadratic characters in condition
(e), since $M\geq1$. A real character $\chi_\theta$ of $E$ is a product of a real character of
conductor dividing $8Q_S'$, which divides $8M$, with $\left(\frac{\cdot}{q}\right)^i\left(\frac{\cdot}{d}\right)^j$;
so (c) also gives the boxes for the real characters of $E$ in condition (e), and the hypothesis of
Lemma~\ref{hf-host-chebyshev} with $B_0=C_0$. That lemma gives condition (f).

Lemma~\ref{hf-norm-primes}, applied to $d$ and $N$, gives primes $g_0,g_1,p,p',q_0$. Since
$p\equiv p'\equiv q_0\equiv1\pmod N$ and $g_i\equiv1\pmod N$ by the congruence on $\pi$, every level
prime $\ell$ is prime to $N$ and $\left(\frac{-d}{\ell}\right)=\left(\frac{\ell}{d}\right)=1$. The
Frobenius of $p$ and $p'$ is $\operatorname{art}(1)=1$ by Proposition~\ref{hf-host-field-properties}(c),
so they split completely in $E$. Corollary~\ref{hf-behrend-set-at-scale}, applied with $L=L_w\leq
C_L\log x$, gives $T_2$.

Lemma~\ref{hf-host-extension} gives a host with level primes $\{g_0,g_1,p,p',q_0\}$,
$\mathfrak c=\mathcal O_K$, the set $T_2$, $u_0=0$, the integer $L_w$, the constant $C_T$ and any top
classes. Lemma~\ref{hf-top-classes} gives top classes $T$ with condition (g) of
Definition~\ref{hf-good-host} and condition (d) of Definition~\ref{hf-host-conditions}. The host
with top classes $T$ satisfies the remaining conditions of Definition~\ref{hf-good-host} by
construction: the splitting bound is $\lfloor z\rfloor\geq\lfloor x^{\varepsilon/(5J)}\rfloor$, and
$J,o\leq C_0\leq C_8$ and $Q_S\leq C_0n\leq C_8n$.
\end{proof}

\section{Hosts on windows of scales}\label{section-hf-windows-hosts}

\begin{definition}\label{hf-hosts-on-windows-def}
Let $\lambda\in\mathbf R$ and let $P$ be a condition on tuples $(C_5,s,R,H,q_0)$ of a real number, a scale,
a rank datum, a host for it and an integer. \emph{Hosts satisfying $P$ exist on windows of exponent}
$\lambda$ if for every $0<\varepsilon<1/52$ there are real numbers $C_T$, $C_8\geq1$, $C_5>0$,
$c_6>0$, $C_b>0$, $x_0$ and a set of scales $(J_W,W)$ with windows of exponent $\lambda$
(Definition~\ref{hf-windows}) such that for every scale $s=(x,\varepsilon,J_W)$ with $x\in W$ and
$x\geq x_0$ there are a rank datum $R$ adapted to $s$ (Definition~\ref{hf-adapted-ranks}), a host $H$ for
$R$ with constant $C_T$ (Definition~\ref{hf-host}) and an integer $q_0$ such that $H$ is good at
$s$ with constants $(C_8,C_5,c_6,C_b)$ and auxiliary prime $q_0$ (Definition~\ref{hf-good-host})
and $P(C_5,s,R,H,q_0)$ holds.
\end{definition}
\leanref{Development/AbelianVarietyTheory/Definitions/GlobalTwisting/ThinScalesIffHostBound.lean}{GlobalTheory.EX.HostExW}

\begin{theorem}\label{hf-hosts-on-windows}
Let $\psi$ be a primitive Dirichlet character that satisfies the weakening of Vandiver's conjecture
with exponent $401$ (Definition~\ref{sha-vandiver-weakening}). Then hosts exist on windows of exponent
$3/2$ (Definition~\ref{hf-hosts-on-windows-def}) that satisfy the following condition at the scale
$s=(x,\varepsilon,J)$, where $t$ is the integer $t$ of the rank datum $R$ and $E$ is the host field
of the host $H$:
\begin{enumerate}
\item the host field satisfies the host conditions for the cyclic level $K_\psi$ of $\psi$
  (Definition~\ref{wl-cyclic-level}) at $(x,\varepsilon,4)$ with top classes $T_{\mathrm{top}}$
  (Definition~\ref{hf-host-conditions});
\item $E$ satisfies the Tate-Shafarevich bound $\bigl(t,\exp(x^{\varepsilon}/2)\bigr)$
  (Definition~\ref{sha-bounded-field}), that is, $\text{Ш}^1_\ell(E,k)$ has finite length for every
  prime $\ell$ and $0<k<t$, and
  \[
   \sum_{\ell\in S}\log\ell\sum_{k=1}^{t-1}\operatorname{len}\text{Ш}^1_\ell(E,k)\leq
   \exp(x^{\varepsilon}/2)
  \]
  for every finite set $S$ of primes.
\end{enumerate}
\end{theorem}
\leanref{Development/AutomorphicArithmetic/CharacterBoxes/SupplyChi.lean}{UpperTheory.Top.supplyChi}
\leanref{Development/AutomorphicArithmetic/ShortNorms/supplyOfAdmOf.lean}{UpperTheory.T3.supplyOfAdmOf}
\leanref{Development/AutomorphicArithmetic/CharacterBoxes/Supply.lean}{UpperTheory.Top.GChi}
\leanref{Development/AutomorphicArithmetic/CharacterBoxes/ExistsPageN.lean}{UpperTheory.Top.VVle}

\begin{proof}
Let $N_\psi$ be the conductor and $o_\psi$ the order of $\psi$, and fix $0<\varepsilon<1/52$.
Proposition~\ref{hf-host-field-exists} gives $c\geq1$ and $e_0$. Proposition~\ref{hf-windows-exist}
with this $c$ gives $J\geq1$ and, for every $y$, a window $[Z,Z^{3/2}]$ with $Z\geq y$; let $W$ be the
set of $x$ for which there are $Y,q,d,e$ with properties (a)--(c) of that proposition. Then $(J,W)$ is a set of scales
with windows of exponent $3/2$. Put $C_0=J+o_\psi+8N_\psi$, and let $C_T,C_8,c_6,C_b$ and $x_0$ be
given by Theorem~\ref{hf-good-host-exists} for $\varepsilon$ and $C_0$. Enlarge $x_0$ so that
Proposition~\ref{hf-adapted-ranks-exist} applies, $x_0>c\,2^{e_0}$, and $Y^{\varepsilon/5}\geq
4N_\psi+5$ whenever $Y\geq x_0^{1/J}$.

Let $s=(x,\varepsilon,J)$ with $x\in W$ and $x\geq x_0$, and take $Y,q,d,e$ with properties
(a)--(c) of Proposition~\ref{hf-windows-exist}. Put $B=\lfloor Y^{\varepsilon/5}\rfloor\geq4N_\psi$. Since
$(q,d)$ is admissible at $Y$, $\left(\frac{-d}{q}\right)=1$ and $\left(\frac{-d}{p}\right)=1$ for odd
$p\leq B$; $q\neq3$ since $\log q>Y^{\varepsilon/5}/2>\log3$; and $\gcd(N_\psi,qd)=1$ by
Definition~\ref{hf-good-pair}(a). Since
$c\,2^e\geq x^{2+8\varepsilon}\geq x>c\,2^{e_0}$, $e\geq e_0$. So
Proposition~\ref{hf-host-field-exists} gives a Galois datum $\Phi$ with group of order $n=c\,2^e$ and index
$o_\psi$, and a host field on $\Phi$ with properties (a)--(e) there. By
Proposition~\ref{hf-windows-exist}(b), $x^{2+8\varepsilon}\leq n\leq2x^{2+8\varepsilon}$, so
Proposition~\ref{hf-adapted-ranks-exist} gives a rank datum $R$ with Galois datum $\Phi$ adapted to $s$. From
$Q_S\leq4N_\psi n$ we get $Q_S\leq C_0n$, and $J,o_\psi,N_\psi\leq C_0$. Theorem~\ref{hf-good-host-exists}
with $M=N_\psi$ gives a host $H$ for $R$ with constant $C_T$ and a prime $q_0$ such that $H$ is good
at $s$ with constants $(C_8,4,c_6,C_b)$, and with top classes satisfying the distribution condition.

For (a), conditions (a), (b) and (c) of Definition~\ref{hf-host-conditions} come from
Proposition~\ref{hf-host-field-exists}: the index is $o_\psi$, which is the integer of $K_\psi$;
property (d) there gives (b) and, since $K_\psi$ has no exceptional character, (c). Condition (d)
comes from Theorem~\ref{hf-good-host-exists}. Condition (e) with $C=4$ is
Lemma~\ref{hf-bounded-characters}: every prime divisor of $Q_S$ divides $2N_\psi$ because $Q_S$
divides $\operatorname{lcm}(2^{e+2},N_\psi)$, and $(q,d)$ is good for $(N_\psi,\varepsilon,Y,x)$.

For (b), Lemma~\ref{hf-kernel-host-field} and property (c) of Proposition~\ref{hf-host-field-exists}
show that $E$ is a host field of $(\psi,e,q,d)$, so Proposition~\ref{hf-windows-exist}(c) gives the
Tate-Shafarevich bound $(\lfloor x^{400}\rfloor,\exp(x^\varepsilon/2))$. Since $t\leq x^{400}$ by
Proposition~\ref{hf-adapted-ranks-exist}, $t\leq\lfloor x^{400}\rfloor$, and the bound holds with
$t$.
\end{proof}
